% concatenated from main-square-class_CME.tex
 \documentclass[11pt]{article}

\usepackage[T1]{fontenc}
\usepackage{lmodern}
\usepackage{algorithm}
\usepackage{algpseudocode}
\usepackage[a4paper,margin=1in]{geometry}
\usepackage{microtype}
\usepackage{amssymb,mathtools}
\usepackage{bm}
\usepackage{booktabs}
\usepackage{xcolor}
\usepackage{pgfplots}
\pgfplotsset{compat=1.18}
\usepackage{enumitem}
\usepackage{amsthm}
\usepackage[hidelinks]{hyperref}
\usepackage[nameinlink,capitalise,noabbrev]{cleveref}

\numberwithin{equation}{section}
\allowdisplaybreaks
\setlist{nosep}
\setlength{\parindent}{1.5em}
\setlength{\parskip}{0pt}

\newtheorem{theorem}{Theorem}[section]
\newtheorem{proposition}[theorem]{Proposition}
\newtheorem{lemma}[theorem]{Lemma}
\newtheorem{corollary}[theorem]{Corollary}
\theoremstyle{definition}
\newtheorem{definition}[theorem]{Definition}
\theoremstyle{remark}
\newtheorem{remark}[theorem]{Remark}




\newcommand{\R}{\mathbb R}
\newcommand{\C}{\mathbb C}
\newcommand{\E}{\mathbb E}
\newcommand{\Var}{\operatorname{Var}}
\newcommand{\SCV}{\operatorname{SCV}}
\newcommand{\Pcal}{\mathcal P}
\newcommand{\Kcal}{\mathcal K}
\newcommand{\Vcal}{\mathcal V}
\newcommand{\Lag}{L}
\newcommand{\dd}{\,\mathrm d}
\newcommand{\e}{\mathrm e}
\newcommand{\ii}{\mathrm i}
\newcommand{\Aop}{\mathcal A}
\newcommand{\lc}{\operatorname{lc}}
\newcommand{\diag}{\operatorname{diag}}
\DeclareMathOperator*{\argmin}{arg\,min}
\newcommand{\fl}{\operatorname{fl}}
\newcommand{\MP}{\mathrm{MP}}
\newcommand{\normw}[1]{\left\lVert #1\right\rVert_{E}}
\newcommand{\ip}[2]{\left\langle #1,#2\right\rangle_{E}}

\title{Least Variability in a Polynomial-Square Class\\ of Rational Kernels}
\author{Maria Laura Battagliola and Oscar Peralta}
\date{}

\begin{document}
\maketitle

\begin{abstract}
We solve an extremal problem for positive randomization kernels that
concentrate a random time as tightly as possible around a deterministic
target, within the class obtained by exponentially damping the square
of a polynomial. The construction, borrowed from the concentrated
matrix-exponential literature, automatically guarantees nonnegativity, gives
the kernel a Laplace transform with a single repeated real pole, and
contains the classical Erlang randomizer as the special case in which the
polynomial is a pure power. The half-polynomial need not be
real-rooted and so may have nonreal conjugate zeros, yet every optimizer is
proved to be real-rooted. After normalizing the kernel to unit mean,
the minimum variance at polynomial degree $m$ turns out to equal the
smallest relative gap between adjacent zeros of the Laguerre polynomial
$L_{m+2}$, with every optimizer obtained by deleting a pair attaining this
minimum. This characterization yields the
minimal variance, the normalized density, and every optimizer in closed form.
The optimizer can be computed from the spectrum of a symmetric tridiagonal
Jacobi matrix. In the large-order limit, the minimal variance decays
quadratically in the order of the kernel, a marked improvement over the
linear decay rate of the Erlang benchmark, and the deleted pair of zeros
localizes at normalized location $2$, with limiting absolute separation
$2\pi$.
\end{abstract}


\paragraph{Keywords.}
Concentrated matrix-exponential distributions; Laguerre polynomials; squared coefficient of variation; positive rational kernels; extremal polynomial problems.



\paragraph{Mathematics Subject Classification (2020).}
Primary 41A20; Secondary 33C45, 60E05.


\section{Introduction}\label{sec:introduction}

Randomization methods replace a fixed evaluation time by a suitably chosen
random horizon. Let $t_0>0$ and let $U$ be a nonnegative random variable
with mean one. For a time-indexed quantity $h\colon[0,\infty)\to\R$, one
then approximates $h(t_0)\approx \E[h(t_0U)]$.
When $U$ is concentrated near one, the randomized time $t_0U$ remains close
to the deterministic target $t_0$. At the same time, an appropriate choice
of the law of $U$ may convert the original fixed-horizon calculation into a
more tractable resolvent, recursive, or matrix-analytic problem. This is the
principle behind maturity randomization and Carr's Canadization procedure in
option valuation \cite{carr1998,bouchardelkarouitouzi2005}, as well as
Erlangization methods for deterministic horizons in ruin theory and related
stochastic models \cite{asmussenavramusabel2002,bladtnielsenperalta2019}.

These methods lead naturally to an extremal problem for the randomizer.
Without structural restrictions the problem is trivial, since the point mass
at $1$ has zero variance. We therefore consider absolutely continuous
randomizers whose densities are nonnegative, whose Laplace transforms are
rational with one prescribed negative real pole location, and whose
time-domain form is an exponentially damped polynomial square. Throughout, a \emph{positive rational kernel} means a nonnegative kernel whose Laplace transform is rational. The
problem is to determine which admissible mean-one kernel has the smallest
variance.

The classical benchmark is the mean-one Erlang randomizer of order $N$, $U\sim \text{Erlang}(N,N)$, whose density and Laplace transform are
\begin{equation}\label{eq:intro-erlang}
\kappa_N^{\mathrm{Er}}(u)
\mathrel{{=}}
\frac{N^N}{(N-1)!}u^{N-1}\e^{-Nu},
\qquad
K_N^{\mathrm{Er}}(s)
\mathrel{{=}}
\left(\frac{N}{N+s}\right)^N.
\end{equation}
Here, $\text{Var}(U) = 1/N$, and its Laplace transform has one negative real pole
of multiplicity $N$. Aldous and Shepp \cite{aldousshepp1987} proved that,
at fixed mean and order, the Erlang distribution has minimum variance among
phase-type distributions. Phase-type distributions form a subclass of
matrix-exponential distributions, which may equivalently be characterized
by nonnegative distributions with rational Laplace transforms
\cite{bladtnielsen2017}. The phase-type subclass has a direct interpretation
as the absorption time of a finite-state Markov jump process, whereas a
general matrix-exponential distribution need not admit such an absorption-time
representation; see also \cite{peralta2023} for a probabilistic interpretation
obtained through exponential tilting. The class studied here is considerably
smaller than the full matrix-exponential class: it has a single real damping
rate and an exponentially damped polynomial-square density, so its transform
has a single repeated real pole. It contains
the Erlang benchmark, but is not restricted to phase-type distributions. We ask
how much the variance can be reduced relative to Erlang within this class while
retaining its repeated-pole structure.

For $a\in\Pcal_m\setminus\{0\}$, where $\Pcal_m$ is the space of real
polynomials of degree at most $m$, and for a damping rate $\lambda>0$, define
\begin{equation}\label{eq:intro-class}
\kappa_{\lambda,a}(u)
\mathrel{{=}}
\frac{\lambda\e^{-\lambda u}a(\lambda u)^2}
{\displaystyle\int_0^\infty \e^{-t}a(t)^2\dd t}.
\end{equation}
By construction, this density is nonnegative. If $\deg a=m$, its Laplace
transform has one distinct pole location, $-\lambda$, of exact multiplicity $N=2m+1$,
so the damping rate fixes the pole location while the polynomial
degree fixes its multiplicity. We choose $\lambda$ so that the kernel
has mean one and minimize its variance over polynomial shapes of exact degree
$m$. The odd-order Erlang density belongs to this class: taking
$a(t)=t^m$ and $\lambda=2m+1$ reduces \eqref{eq:intro-class} to
\eqref{eq:intro-erlang} with $N=2m+1$.

No real-rootedness assumption is imposed on $a$. Since $a$ has real
coefficients, any nonreal zeros occur in conjugate pairs, and these are
zeros of the polynomial factor rather than poles of the Laplace transform,
so the transform still has the single real pole location $-\lambda$ even
when $a$ has nonreal zeros.

Expanding $a^2$ shows that the density $\kappa_{\lambda,a}$ in
\eqref{eq:intro-class} is a signed linear combination of common-rate Erlang
densities of orders $1$ through $N$, an expansion recorded precisely in
Remark~\ref{aw:rem:confluent}. We refer to $U$, or equivalently to its law, as the \emph{randomizer}. Any quantity obtained from the randomizer by a linear
operation, such as the expectation $\E[h(t_0U)]$ or the Laplace transform
$K(s)$, therefore decomposes into the same operation applied to each Erlang
component, recombined with the signed weights from the expansion.

Variance is the natural concentration criterion from both the randomization
and transform viewpoints. We show that, under mild regularity conditions on $h$, the
error produced by replacing $t_0$ with $t_0U$ is controlled by
$\Var(U)$. 
Similarly, if
$K(s)=\E[\e^{-sU}]$, then matching the mass and mean of the degenerate randomizer
makes the first discrepancy between $K(s)$ and $\e^{-s}$ near the
origin depend on the variance. Smaller variance therefore means both a more
concentrated randomized horizon and closer local agreement with the
deterministic Laplace transform. The latter interpretation is developed in
\cref{sec:positive-kernels}. The optimization considered here is,
nevertheless, not a problem of best uniform, weighted-uniform, or
integral-norm rational approximation of $\e^{-s}$ on a nontrivial interval.
It is a constrained moment extremal problem within the class
\eqref{eq:intro-class}.

Beyond the probabilistic motivation described above, the problem connects with several established theories. One such connection is with the
theory of concentrated matrix-exponential (CME) distributions. The concentrated matrix-exponential line of research was initiated by Horv\'ath, S\'af\'ar, Telek, and Z\'amb\'o \cite{horvathsafartelekzambo2016}, and the connection between
concentrated nonnegative matrix-exponential functions and positive,
non-overshooting randomization kernels was developed by Horv\'ath,
Talyig\'as, and Telek \cite{horvathtalyigastelek2018}. Stable high-order
constructions and refinements of the associated nonlinear optimization were
given in
\cite{horvathhorvathtelek2020,almousatelek2021}, while the corresponding
numerical inverse Laplace-transform method and its implementation were
studied in \cite{horvathetal2020nilt,almousaetal2022cme}. M\'esz\'aros and
Telek subsequently formulated the real-eigenvalue branch, denoted CME-R,
and developed numerical procedures for finding highly concentrated
real-eigenvalue matrix-exponential distributions \cite{meszarostelek2022}.
The full CME-R problem is broader than the one treated here because it can
involve several distinct real pole locations. Within their numerical study,
a tractable repeated-pole family is parameterized by real zeros,
\[
 a(z)=c\prod_{r=1}^m(z-\tau_r),\qquad \tau_r\in\R,
\]
so that the density is proportional to
$\e^{-\lambda u}\prod_{r=1}^m(\lambda u-\tau_r)^2$. The class
\eqref{eq:intro-class} keeps the same single repeated real pole and
square-polynomial structure but drops the real-rootedness restriction on
$a$. Thus, it is a subclass of full CME-R, but an enlargement of this
real-rooted repeated-pole family (strict for $m\geq2$). Nonreal conjugate
zeros of $a$ do not introduce complex ME eigenvalues because they are zeros
of the time-domain polynomial factor rather than transform poles.

M\'esz\'aros and Telek \cite{meszarostelek2022} reported numerical candidates in the
real-rooted repeated-pole family and observed an approximate $m^{-1.85}$
decay at orders $N=2m+1$. The theorem below optimizes over the larger polynomial-square
family \eqref{eq:intro-class}, which allows the polynomial factor $a$ to have nonreal conjugate roots, and proves that every optimizer is nevertheless
real-rooted. Consequently, enlarging the search space to allow nonreal zeros
does not lower the optimum, since the minimum over our class equals the minimum
over their real-rooted repeated-pole family. What remains outside the scope
of the theorem is the genuinely broader CME-R problem with more general
real-pole architectures.

Broader CME and CME-R parameterizations report smaller finite-order SCVs at
some orders, but these values are obtained by order-by-order numerical
searches and are not accompanied by global optimality guarantees for the
full classes. The result below instead certifies every optimizer within the
repeated-real-pole square-polynomial subclass.

The same problem also has an approximation-theoretic interpretation.
Rational approximation with prescribed or concentrated poles classically
asks for best uniform or weighted-uniform approximation of the exponential
by general rational functions. Fixed-pole degree estimates
\cite{anderssonganelius1977}, concentrated-pole approximation of the decaying exponential on the positive half-line \cite{andersson1981}, and related real-pole constructions
\cite{borwein1983} belong to
this tradition. Orthogonal rational functions provide a broader setting
for prescribed-pole questions \cite{bultheeletal1999}, while Pad\'e
approximation and quadrature are closely connected to rational approximation
of the exponential more generally
\cite{bakergravesmorris,prevostrivoal2007}. The problem considered here is an extremal polynomial problem, in the broad tradition of relations between optimal measures and extremal polynomial growth \cite{boslevenbergortega2021}. Its distinctive feature is that positivity, normalization, a fixed pole architecture, and a moment-concentration objective are imposed simultaneously. Through the representation \eqref{eq:intro-class}, these conditions become constraints on the polynomial factor, while variance minimization supplies the extremal criterion.


Our solution proceeds by converting the variance minimization into a
sequence of polynomial and spectral reductions. Scale invariance and a
change of variables first transform the nonlinear moment ratio into a
Rayleigh-type quotient centered at a fixed point. The stationarity
conditions force the extremizing polynomial to satisfy a system of
Laguerre orthogonality relations. These relations yield a rigidity result that reduces the optimizer, after removing one additional degree of freedom, to a single Laguerre mode. Applying Gauss--Laguerre quadrature then shows that the two zeros removed from this mode must be adjacent. Finally, strict monotonicity of the minimum logarithmic Laguerre zero spacing across consecutive degrees selects the exact polynomial degree. Since Laguerre zeros are the eigenvalues of a
symmetric tridiagonal Jacobi matrix, the same argument produces an
exhaustive construction.

To state the result, let $\Lag_n$ denote the ordinary Laguerre polynomial,
with zeros $0<x_{1,n}<\cdots<x_{n,n}$,
and define, for adjacent zeros,
\begin{equation}
\label{eq:deltas_intro}
 {\gamma_{j,n}}
=
\frac{(x_{j+1,n}-x_{j,n})^2}
{4x_{j,n}x_{j+1,n}},
\qquad
\Delta_n
=
\min_{1\leq j<n}{\gamma_{j,n}}. 
\end{equation}
Our main theorem shows that the minimum variance at polynomial degree $m$ is
$\mathfrak v_m:=\Delta_{m+2}$, and that every minimizing polynomial is obtained by deleting
a minimizing adjacent pair of zeros from $\Lag_{m+2}$. In rational order
$N=2m+1$, the optimal variance satisfies
$$
\mathfrak v_m\sim\frac{\pi^2}{N^2}.
$$
The Erlang benchmark has variance $1/N$, so enlarging the admissible class
from phase-type Erlang randomizers to the positive rational kernels considered
here changes the attainable concentration rate from order $N^{-1}$ to order
$N^{-2}$. The asymptotic result follows from a nonasymptotic logarithmic-gap
bound and the Marchenko--Pastur limit for normalized Laguerre zeros. The
same analysis localizes every minimizing pair at the normalized point $2$,
and shows that the absolute separation of the deleted zeros tends to
$2\pi$.

For comparison, the optimization-free construction of Battagliola and
Peralta \cite{battagliolaperalta2026} has squared coefficient of variation of
order $\log^3(N)/N^2$. The present $N^{-2}$ rate is therefore sharper, and its proof uses Laguerre geometry and quadrature
rather than an explicit trigonometric construction.

The result provides complementary advantages in each of the connected
theories. For the square-polynomial with a single repeated real pole part of the CME-R problem, it replaces order-by-order nonlinear
optimization by an exact characterization of every optimizer, together with
a Jacobi matrix construction and a sharp quadratic asymptotic law.
From the approximation-theoretic viewpoint, it gives an explicit solution
to a positivity-constrained moment extremal problem and identifies the
extremizer through an adjacent-zero deletion rule for Laguerre polynomials.
From the randomization viewpoint, it produces repeated-pole kernels that are
asymptotically much more concentrated than the Erlang randomizer. Our optimality result is restricted to kernels of the form \eqref{eq:intro-class}, in which nonnegativity is enforced by the square of a single real polynomial and the Laplace transform has one negative real pole location of prescribed multiplicity. We do not
claim optimality over the full cone of polynomials nonnegative on
$[0,\infty)$, over kernels with several distinct or complex poles, over
signed kernels, or with respect to a global transform norm. The result is
also complementary to the construction in
\cite{battagliolaperalta2026}, which uses powered Fej\'er kernels to obtain
explicit concentrated matrix-exponential laws with split
common-damping complex poles.

The paper is organized as follows.
\Cref{sec:foundations} develops the randomization identities and the
square-polynomial repeated-pole class.
\Cref{sec:variational} reduces the problem to a fixed-center Rayleigh
quotient.
\Cref{sec:laguerre-rigidity} establishes the one-mode Laguerre rigidity
result.
\Cref{sec:adjacency} proves the adjacent-zero criterion, determines the exact
degree, and states the finite-order optimizer.
\Cref{sec:computation} develops the Jacobi spectral construction.
\Cref{sec:asymptotics} proves the sharp large-order asymptotics and the
localization of the deleted zeros.
Finally, \cref{sec:discussion} compares the randomization,
matrix-exponential, and approximation-theoretic interpretations.








\section{Positive randomization and the square-polynomial repeated-pole class}\label{sec:foundations}

This section develops the randomization framework and introduces the admissible kernel class on which the subsequent extremal analysis is built. Specifically, \Cref{sec:positive-kernels} formulates positive randomization as an
operator on test functions, and proves the exact time-domain and
transform-domain identities that make the variance the appropriate optimization
criterion. \Cref{sec:repeated-class} then introduces the
square-polynomial repeated-pole class itself, together with its rational
Laplace transform structure and the precise optimization problem studied
in the remainder of the paper.

\subsection{Positive randomization and exact concentration error}
\label{sec:positive-kernels}

Let $U$ be a nonnegative random variable with density $\kappa$, and define
\begin{equation}\label{rand:eq:operator}
  \Aop_\kappa h(t)
  :=\int_0^\infty h(tu)\kappa(u)\dd u
  =\E[h(tU)].
\end{equation}

The deterministic evaluation operator is recovered in the degenerate case $U=1$ almost surely, and the next result quantifies the error introduced by replacing this degenerate randomizer with a nondegenerate randomizer. 
Unlike a signed approximation formula, thanks to its expectation representation, the operator \eqref{rand:eq:operator} enjoys favorable properties by construction. First, when $h$ is bounded,
$\inf_{t\geq0}h(t)\leq\Aop_\kappa h(t)\leq\sup_{t\geq0}h(t)$. Second, when $h$ is
increasing, $t\mapsto\Aop_\kappa h(t)$ is increasing, and the same holds for
decreasing $h$. Third, when $h\geq0$, so is $\Aop_\kappa h$.


The natural centering condition is $\E[U]=1$. This entails that
\begin{align*}
  \Var(U)=\E[(U-1)^2],
\end{align*}
is also the squared coefficient of variation, which classically is the most immediate measure
of concentration about the deterministic target.

\begin{proposition}
\label{rand:prop:error-bound}

Suppose that $\E[U]=1$ and $\Var(U)<\infty$. Then, for every $t>0$,
\begin{equation*}
 \sup_{\substack{h\in C^2([0,\infty))\\ \lVert h''\rVert_\infty\leq1}}
 \left|\E[h(tU)]-h(t)\right|
 =\frac{t^2}{2}\Var(U).
\end{equation*}
Here $C^2([0,\infty))$ denotes the space of functions on
$[0,\infty)$ with two continuous derivatives, and
$\lVert h''\rVert_\infty=\sup_{x\geq0}|h''(x)|$ is the supremum norm of the
second derivative.
\end{proposition}


\begin{proof}

Taylor's formula with integral remainder gives, for every admissible $h$,
\begin{align*}
 h(tu)
 &=h(t)+t h'(t)(u-1)\\
 &\quad+t^2(u-1)^2
 \int_0^1(1-r)h''\left(t+rt(u-1)\right)\dd r.
\end{align*}
After taking expectations, the linear term vanishes and the remainder is
bounded in absolute value by $t^2\Var(U)/2$. Equality is attained by
$h_0(t)=t^2/2$, for which $\lVert h_0''\rVert_\infty=1$ and
\[
 \E[h_0(tU)]-h_0(t)=\frac{t^2}{2}\bigl(\E[U^2]-1\bigr)
 =\frac{t^2}{2}\Var(U).
\]

\end{proof}


The variance also has a direct local transform-domain meaning. The Laplace
transform of the randomizer is
\begin{align*}
  K(s)=\E[\e^{-sU}]
      =\int_0^\infty \e^{-su}\kappa(u)\dd u.
\end{align*}
Since the degenerate randomizer $U=1$ has Laplace transform $\e^{-s}$, we compare $K(s)$ with $\e^{-s}$ locally around $s=0$. In particular, if $\E[U]=1$ and $\E[U^2]<\infty$, then
\begin{align*}
  \mathopen{{\lvert}}K(s)-\e^{-s}\mathclose{{\rvert}}
  =\frac{\Var(U)}2s^2+o(s^2),
  \qquad s\downarrow0.
\end{align*}
This follows from the second-order Taylor expansions of $K(s)$ and $\e^{-s}$
at the origin.
Hence, minimizing variance minimizes
the first unmatched Taylor coefficient after mass and mean matching. 

The mean-one Erlang randomizer provides the classical benchmark. Its phase-type representation as a sum of exponential stages makes the randomized horizon accessible through finite-state Markovian recursions, as exploited in Canadization for option pricing and in Erlangization for finite-horizon ruin probabilities \cite{carr1998,asmussenavramusabel2002}, but its variance at order $N$ is only $1/N$. The class studied here retains the positive rational-transform structure and the single repeated real pole while allowing matrix-exponential laws that need not be phase-type. This broader class is relevant in state-space-sensitive constructions such as ME-fication, where a more concentrated randomizer at the same matrix order can reduce the dimension of the auxiliary systems used to recover transient and first-passage quantities \cite{akargursoyhorvathtelek2021}. Such matrix-exponential randomizers do not, however, automatically inherit every probabilistic recursion available for Erlang randomizers.



The randomization operator introduced above also gives an integral interpretation of the Abate--Whitt inversion framework, which we briefly recall before relating it to the repeated-pole setting. For a Laplace transform $\widehat h(s)=\int_0^\infty \e^{-st}h(t)\dd t$,
the Abate--Whitt framework approximates $h(T)$ by a finite linear combination of values of $\widehat h$ at the scaled points $\beta_k/T$, where the complex node parameters $\beta_k$ and weights $\eta_k$ depend on the approximation order but not on $\widehat h$ or $T$ \cite{abatewhitt2006}. More precisely, if an approximate-identity kernel has the exponential representation $f_n(u)=\operatorname{Re}\mathopen{{\bigl(}}\sum_{k=1}^n \eta_k\e^{-\beta_k u}\mathclose{{\bigr)}}$,
then a change of variables gives
\begin{align*}
 \int_0^\infty h(Tu)f_n(u)\dd u
 =\frac1T\operatorname{Re}\mathopen{{\biggl(}}\sum_{k=1}^n
   \eta_k\widehat h\!\left(\frac{\beta_k}{T}\right)\mathclose{{\biggr)}}.
\end{align*}
This is the corresponding Abate--Whitt inversion formula. Whenever $f_n$ is nonnegative and has unit mass, this evaluation is a positive randomization, as it preserves the range and monotonicity of $h$ and produces neither positive nor negative overshoot. The CME method lies in this nonnegative part of the Abate--Whitt framework and selects kernels of this form that are highly concentrated around one \cite{horvathtalyigastelek2018,horvathetal2020nilt}.




The repeated-pole family studied here provides a repeated-node counterpart of the
Abate--Whitt representation. Whereas a sum of exponential terms leads to
transform evaluations at distinct scaled nodes, the polynomial modulation in
\eqref{eq:intro-class} produces derivative evaluations at a single point.
Whenever $h$ is regular enough that $\widehat h$ can be
differentiated under the Laplace integral through order $2m$ at
$s=\lambda/T$, a mild condition satisfied for instance by every $h$ of at
most polynomial growth on $[0,\infty)$, writing $a(z)^2=\sum_{j=0}^{2m}c_jz^j$
and differentiating under the Laplace integral gives
\begin{align*}
\Aop_{\kappa_{\lambda,a}}h(T)
=\frac{\lambda}
{T\displaystyle\int_0^\infty \e^{-z}a(z)^2\dd z}
\sum_{j=0}^{2m}c_j
\left(-\frac{\lambda}{T}\right)^j
\widehat h^{(j)}\left(\frac{\lambda}{T}\right).
\end{align*}
Hence, a pole of multiplicity $N=2m+1$ corresponds to evaluating
$\widehat h$ and its first $2m$ derivatives at the single point
$\lambda/T$, the repeated-node analogue of evaluating $\widehat h$ at
several distinct nodes. From this viewpoint, the present extremal problem selects the
nonnegative mean-one kernel in this repeated-pole class with minimum
variance.




\subsection{The square-polynomial repeated-pole class}\label{sec:repeated-class}

We now construct the polynomial-square family used in the definition of the kernel \eqref{eq:intro-class} precisely. Fix $m\ge0$ and let
\[
    \Pcal_m
    =\left\{a(z)=\sum_{r=0}^{m}a_rz^r:a_r\in\R\right\}.
\]
For $0\ne a\in\Pcal_m$ and $k=0,1,2$, define
\begin{equation}\label{aw:eq:Qk}
    Q_k(a)=\int_0^\infty z^k\e^{-z}a(z)^2\dd z.
\end{equation}
Then, for a damping rate $\lambda>0$, set
\begin{equation}\label{aw:eq:kappa-lambda-a}
    \kappa_{\lambda,a}(u)
    =\frac{\lambda\e^{-\lambda u}a(\lambda u)^2}{Q_0(a)},
    \qquad u\ge0.
\end{equation}
The denominator $Q_0(a)$, together with the change of variables $z=\lambda u$, ensures that kernel \eqref{aw:eq:kappa-lambda-a} integrates to one, and nonnegativity is ensured because the polynomial factor is squared.
The overall scale of $a$ is irrelevant, since $\kappa_{\lambda,ca}=\kappa_{\lambda,a}$
for every $c\ne0$, since the resulting factor cancels between the numerator and denominator. Thus, the shape of the kernel depends only on the relative coefficients of $a$. In other words, the relevant parameter is the projective direction of the coefficient vector of $a$, and one coefficient normalization may be imposed without loss of generality.



The rationality of the transform, and the exact multiplicity of its pole,
follow from the same square-polynomial structure. Writing $a(z)^2=\sum_{r=0}^{2m}d_rz^r$,
the Laplace transform of \eqref{aw:eq:kappa-lambda-a} is
\begin{align}
    K_{\lambda,a}(s)
    &=\int_0^\infty\e^{-su}\kappa_{\lambda,a}(u)\dd u\notag\\
    &=\frac1{Q_0(a)}
      \int_0^\infty
      \exp\!\left[-\left(1+\frac{s}{\lambda}\right)z\right]
      a(z)^2\dd z\notag\\
    &=\frac1{Q_0(a)}
      \sum_{r=0}^{2m}d_rr!
      \left(1+\frac{s}{\lambda}\right)^{-(r+1)}\notag\\
    &=\frac1{Q_0(a)}
      \sum_{r=0}^{2m}d_rr!
      \left(\frac{\lambda}{s+\lambda}\right)^{r+1}.
      \label{aw:eq:K-rational}
\end{align}
Therefore, $K_{\lambda,a}$ is rational and has only one possible pole at
$s=-\lambda$. If $\deg a=m$, then $d_{2m}>0$, and the pole has exact
multiplicity $N=2m+1$. This denominator degree is the order used in the rational and
repeated-pole statements below.


\begin{remark}
\label{aw:rem:confluent}
The single damping factor $\e^{-\lambda u}$ determines the unique pole
location $-\lambda$. The polynomial degree determines the multiplicity: since
$a^2$ has degree $2m$, the largest denominator power in
\eqref{aw:eq:K-rational} is $(s+\lambda)^{2m+1}$. If
\[
 g_{r+1,\lambda}(u)
 =\frac{\lambda^{r+1}u^r\e^{-\lambda u}}{r!},
 \qquad r\geq0,
\]
denotes the Erlang density of order $r+1$ and rate $\lambda$, then
\begin{equation*}
 \kappa_{\lambda,a}(u)
 =\sum_{r=0}^{2m}w_r g_{r+1,\lambda}(u),
 \qquad
 w_r=\frac{d_r r!}{Q_0(a)}.
\end{equation*}
The weights sum to one but may have either sign. Nonnegativity of the total
density follows from the square $a(\lambda u)^2$. Any calculation linear in
the randomizing law can therefore be performed for each Erlang component and
recombined with the signed weights.
\end{remark}

Formula \eqref{aw:eq:K-rational} also explains the matrix-exponential
terminology. It is a proper real rational Laplace transform with all poles
in the open left half-plane, and hence admits a finite-dimensional real
matrix-exponential realization. Because all poles coalesce at the one real
location $-\lambda$, the present class is the repeated-real-pole
square-polynomial subclass of the full CME-R problem discussed in the
introduction, and is strictly larger than the real-rooted repeated-pole
parametrization used in the numerical study of \cite{meszarostelek2022}.
The restriction to one repeated pole and one polynomial square, however,
remains part of every optimality statement below.

Finally, with the class of kernels and its rational structure in hand, we can now
state the optimization problem precisely. If $T_a$ has density $\kappa_{1,a}$
and $U_{\lambda,a}=T_a/\lambda$, then
\begin{align}
\label{num:eq:mean}
    \E[U_{\lambda,a}]
    &=\frac{Q_1(a)}{\lambda Q_0(a)},\\
 \label{num:eq:second}
    \E[U_{\lambda,a}^2]
    &=\frac{Q_2(a)}{\lambda^2Q_0(a)},
\end{align}
and there is a unique scale that centers the kernel at one, namely
\begin{align*}
    \lambda(a)=\frac{Q_1(a)}{Q_0(a)}.
\end{align*}
Moreover, the squared coefficient of variation of $T_a$ is
\begin{equation}
\label{eq:scv}
  \SCV(T_a):=\frac{\Var(T_a)}{\E[T_a]^2}\mathpunct{{.}}  
\end{equation}
Notice that, since $T_a$ is a nondegenerate absolutely continuous random variable, $\Var(T_a)>0$ and hence $\SCV(T_a)>0$. This allows us to write the target function to optimize 
\begin{equation}
\label{eq:obj_scv}
    \Vcal(a)
    :=\Var(U_{\lambda(a),a})
    =\frac{Q_0(a)Q_2(a)}{Q_1(a)^2}-1
    =\SCV(T_a).
\end{equation}
The least-variable repeated-pole randomization problem is therefore the
following.

\begin{definition}[Repeated-pole randomization problem]
\label{aw:def}
For each $m\geq0$, define the optimal variance among the mean-one
randomizers of polynomial degree $m$ by
\begin{align*}
\mathfrak v_m
:=
\inf_{\substack{a\in\Pcal_m\ \deg a=m}}
\Vcal(a).
\end{align*}
\end{definition}
That is, $\mathfrak v_m$ is the infimum of the variances over
mean-one kernels $\kappa_{\lambda(a),a}$ whose polynomial factor has exact
degree $m$. Equivalently, its Laplace transform has one real pole of exact
multiplicity $2m+1$.
At the endpoint $m=0$, the polynomial $a$ is constant, the mean-one
normalization gives $\lambda=1$, and the resulting kernel is the unit-rate
exponential density. Hence 
\begin{equation}
\label{eq:case_m0}
    \mathfrak v_0=1=\Delta_2,
\end{equation}
according to the notation introduced in \eqref{eq:deltas_intro}.
The structural arguments below are written for $m\geq1$, while the main
theorem includes $m=0$ through this direct calculation.



\section{Variational characterization}\label{sec:variational}
In this section and through the degree-selection theorem
(\Cref{adj:thm:final}), assume $m\ge1$.
The endpoint $m=0$ corresponds to the exponential case, and admits a single
squared coefficient of variation equal to $1$ as shown in \eqref{eq:case_m0}, so it is trivial.

For a nonzero polynomial $a\in\Pcal_m$, the quantities $Q_k(a)$ from
\eqref{aw:eq:Qk} are finite and strictly positive, and the associated
unit-damping density is $\kappa_{1,a}$. As in the previous section, we denote by $T_a$ a random variable with this density. 


We first consider the relaxed problem obtained by replacing the exact-degree
condition $\deg a=m$ with $\deg a\leq m$. After identifying nonzero scalar
multiples, the exact-degree class is an open subset of the projective
parameter space because its leading coefficient may tend to zero along a
convergent sequence, and its closure consists of all nonzero polynomials of
degree at most $m$. This closure is compact, so the continuous objective
attains its minimum there. We therefore begin with
\begin{equation*}
\inf_{0\ne a\in\Pcal_m}\Vcal(a),
\end{equation*}
where $\Vcal(a):=\SCV(T_a)$ is defined in \eqref{eq:scv}. We later determine whether the minimizer has degree $m-1$ or $m$ and show, through a strict comparison of these two outcomes, that every minimizer has exact degree $m$.



The polynomial $a$ itself
is a half-polynomial of the density, defined only up to a nonzero scalar,
so before existence can be addressed we must first record how the moment
functionals behave under this rescaling.
For $a\ne0$ and $c\ne0$, substituting $ca$ into
\eqref{aw:eq:Qk} and \eqref{aw:eq:kappa-lambda-a} with fixed $\lambda=1$ gives
\begin{equation*}
 Q_k(ca)=c^2Q_k(a), \qquad \kappa_{1,ca}=\kappa_{1,a},
 \quad k=0,1,2.
\end{equation*}
Every moment objective therefore depends on $a$ only through the
one-dimensional subspace it spans, or equivalently through its projective
class, and is unchanged when $a$ is replaced by any nonzero scalar multiple,
so the objective 
descends directly to the real
projective space $\mathbb{RP}^m$ of nonzero coefficient vectors modulo
nonzero scalar multiplication. This space is compact, and no normalization
of a representative is required.
If $a(t)=\sum_{j=0}^m a_jt^j$, then
\begin{equation*}
Q_k(a)=\sum_{i=0}^m\sum_{j=0}^m(i+j+k)!{\,}a_i a_j,
\qquad k=0,1,2.
\end{equation*}
Hence, each $Q_k$ is continuous in the coefficients of $a$. Although the
individual functionals $Q_k$ depend on the chosen representative of the
projective class, their scale factors cancel in the moment ratios defining
$\Vcal$. Consequently, $\Vcal$ defines a continuous function on the
compact projective parameter space.
It follows that $\Vcal$ attains its minimum over the degree-at-most-$m$ class. This does not yet ensure that a minimizer has exact degree $m$, and lower-degree minimizers are ruled out later by a separate comparison argument.




Existence settled, we turn to a reformulation of the objective as a relative quadratic loss, which is the form the variational argument actually needs. 
 As in \eqref{eq:obj_scv}, we know that the function to be minimized corresponds to $\SCV(T_a)=\frac{Q_0(a)Q_2(a)}{Q_1(a)^2}-1$.
For the variational calculation, it is convenient to use
\begin{align*}
    \rho(a)
    =\frac{\Var_a(T_a)}{\E_a[T_a^2]}
    =1-\frac{Q_1(a)^2}{Q_0(a)Q_2(a)}.
\end{align*}
The relation between \(\rho\) and the squared coefficient of variation is
\begin{align*}
    \rho(a)=\frac{\SCV(T_a)}{1+\SCV(T_a)}\in (0,1),
    \qquad
    \SCV(T_a)=\frac{\rho(a)}{1-\rho(a)} \in (0,\infty).
\end{align*}
The transformation between $\SCV(T_a)$ and $\rho(a)$ maps $(0,\infty)$ strictly increasingly into $(0,1)$. Therefore,
minimizing \(\rho\) is equivalent to minimizing \(\SCV\).  Let
\begin{align*}
    \rho_*=\min_{0\ne a\in\Pcal_m}\rho(a),
\end{align*}
and choose a minimizing representative \(a_*\in\Pcal_m\setminus\{0\}\), which
by scale invariance is determined only up to multiplication by a nonzero
constant. The transformed objective \(\rho(a)\) admits a useful
representation as an optimized relative quadratic loss. In particular, it introduces
an auxiliary center playing the role of the mean-one damping rate once the
polynomial is fixed.

\begin{proposition}\label{var:prop:relative-loss}
Let \(T\) be a positive random variable with \(0<\E[T]<\infty\) and \(\E[T^2]<\infty\). Then
\begin{equation*}
    \frac{\Var(T)}{\E[T^2]}
    =\min_{\mu>0}\frac{\E[(T-\mu)^2]}{\mu^2}.
\end{equation*}
The minimizer is unique and equals 
\begin{align*}
\frac{\E[T^2]}{\E[T]}.
\end{align*}
\end{proposition}
\begin{proof}
Write \(m_1=\E[T]\) and \(m_2=\E[T^2]\). For \(\mu>0\), expanding the square gives
\begin{equation}\label{var:eq:Fmu-expansion}
    G(\mu):=\frac{\E[(T-\mu)^2]}{\mu^2}
    =\frac{m_2}{\mu^2}-\frac{2m_1}{\mu}+1,
\end{equation}
so
\[
    G'(\mu)=\frac{2(m_1\mu-m_2)}{\mu^3},
\]
which is negative for \(0<\mu<m_2/m_1\) and positive beyond, giving the unique global minimizer at \(m_2/m_1\). Substituting back into \eqref{var:eq:Fmu-expansion},
\[
    G\left(\frac{m_2}{m_1}\right)
    =1-\frac{m_1^2}{m_2}
    =\frac{m_2-m_1^2}{m_2}
    =\frac{\Var(T)}{\E[T^2]},
\]
which proves both assertions.
\end{proof}

Applying Proposition~\ref{var:prop:relative-loss} to \(T_a\) gives
\begin{align*}
    \mu(a)
    =\frac{\E_a[T_a^2]}{\E_a[T_a]}
    =\frac{Q_2(a)}{Q_1(a)},
\end{align*}
where $\mu(a)$ denotes the unique minimizer identified in Proposition~\ref{var:prop:relative-loss}, now viewed as a function of $a$. We write $\mu_*$ in place of $\mu(a_*)$ once the optimal polynomial $a_*$ is fixed, and more generally drop the argument whenever the polynomial in question is clear from context, but the dependence on $a$ remains implicit throughout.
To derive a first-order optimality condition in the polynomial coefficients, we introduce the center as an auxiliary variable and consider the joint quotient
\begin{equation}\label{var:eq:J-def}
\mathcal J(a,\mu) =
\frac{
\displaystyle\int_0^\infty (t-\mu)^2a(t)^2e^{-t},dt
}{
\displaystyle\mu^2\int_0^\infty a(t)^2e^{-t},dt
},
\qquad
0\ne a\in\Pcal_m,\quad \mu>0.
\end{equation}
For each fixed $a$, minimizing over the center recovers the original objective,
$$
\rho(a)=\min_{\mu>0}\mathcal J(a,\mu),
\qquad
\mu(a)=\frac{Q_2(a)}{Q_1(a)}.
$$
Consequently,
$$
\rho_*
=
\min_{\substack{0\ne a\in\Pcal_m\ \mu>0}}
\mathcal J(a,\mu),
$$
and if $a_*$ minimizes $\rho$, then
$$
\mu_*=\frac{Q_2(a_*)}{Q_1(a_*)},
$$
makes $(a_*,\mu_*)$ a joint minimizer of $\mathcal J$. We now hold this optimal center $\mu_*$ fixed and vary only the polynomial. Joint optimality then implies that $a_*$ minimizes
$$
\Lambda_{\mu_*}(a)
=
\frac{
\displaystyle\int_0^\infty (t-\mu_*)^2a(t)^2e^{-t},dt
}{
\displaystyle\int_0^\infty a(t)^2e^{-t},dt
},
\qquad
0\ne a\in\Pcal_m,
$$
because
$$
\mathcal J(a,\mu_*)=\frac{\Lambda_{\mu_*}(a)}{\mu_*^2}.
$$
The first-order optimality condition for this Rayleigh quotient yields
the Euler equation below. Although $\mu_*$
 is not yet known explicitly, it is treated as fixed when taking variations in the polynomial $a$.



\begin{theorem}[First variation in polynomial directions]\label{var:thm:first-variation}
Let \((a_*,\mu_*)\) be a joint minimizer of \(\mathcal{J}\) in \eqref{var:eq:J-def}, and let \(\rho_*=\mathcal{J}(a_*,\mu_*)\). Define
\begin{equation}\label{var:eq:eta-def}
    \eta_*
    :=\rho_*\mu_*^2
    =\min_{0\ne a\in\Pcal_m}\Lambda_{\mu_*}(a),
\end{equation}
with $a_*$ realizing this minimum, and set
\begin{equation}\label{var:eq:R-def}
    R_*(t)=(t-\mu_*)^2-\eta_*.
\end{equation}
Then
\begin{equation}\label{var:eq:Euler-integral}
    \int_0^\infty R_*(t)a_*(t)b(t)\e^{-t}\,dt=0
    \qquad
    \text{for every }b\in\Pcal_m.
\end{equation}
\end{theorem}

\begin{proof}
For \(b\in\Pcal_m\), consider the path \(a_\varepsilon=a_*+\varepsilon b\), admissible since \(a_*\ne0\) keeps \(a_\varepsilon\) nonzero for small \(\varepsilon\). Holding \(\mu=\mu_*\) fixed and writing
\begin{align*}
    N(\varepsilon)
    &=\int_0^\infty
      (t-\mu_*)^2a_\varepsilon(t)^2\e^{-t}\,dt,
      &
    D(\varepsilon)
    &=\mu_*^2\int_0^\infty
      a_\varepsilon(t)^2\e^{-t}\,dt,
\end{align*}
we have \(\mathcal{J}(a_\varepsilon,\mu_*)=N(\varepsilon)/D(\varepsilon)\). Since \((a_*,\mu_*)\) is a joint minimizer, this ratio has a local minimum at \(\varepsilon=0\), so
\begin{equation}\label{var:eq:stationarity-epsilon}
    N'(0)D(0)-N(0)D'(0)=0.
\end{equation}
Expanding \(a_\varepsilon^2=a_*^2+2\varepsilon a_*b+\varepsilon^2b^2\) gives
\begin{align}
    N'(0)
    &=2\int_0^\infty
      (t-\mu_*)^2a_*(t)b(t)\e^{-t}\,dt,
      \label{var:eq:N-prime}\\
    D'(0)
    &=2\mu_*^2\int_0^\infty
      a_*(t)b(t)\e^{-t}\,dt.
      \label{var:eq:D-prime}
\end{align}
Since \(\rho_*=N(0)/D(0)\), \eqref{var:eq:stationarity-epsilon} reduces to \(N'(0)-\rho_*D'(0)=0\). Inserting \eqref{var:eq:N-prime}--\eqref{var:eq:D-prime}, dividing by two, and combining the integrals gives
\begin{align*}
    0
    &=\int_0^\infty
      \left[(t-\mu_*)^2-\rho_*\mu_*^2\right]
      a_*(t)b(t)\e^{-t}\,dt
    =\int_0^\infty
      R_*(t)a_*(t)b(t)\e^{-t}\,dt,
\end{align*}
which proves \eqref{var:eq:Euler-integral}.
\end{proof}

\begin{remark}\label{var:rem:mu-fixed} 
Since $(a_*,\mu_*)$ is a joint minimizer,
stationarity holds along paths that vary $a$ while keeping $\mu=\mu_*$.
Equivalently, after substituting the minimizing center $\mu(a)$, the chain-rule
term involving $D\mu(a_*)$ vanishes because
$\partial_\mu\mathcal J(a_*,\mu_*)=0$. The symbols $\rho_*$, $\mu_*$, and
$\eta_*$ remain written in terms of $a_*$ only because the optimizer is
initially unknown. 
\end{remark}

\begin{remark}
Write the polynomial in the monomial basis as
$$
a(t)=\sum_{i=0}^m a_i t^i,
\qquad
\bm a=(a_0,\ldots,a_m)^\mathsf{T},
$$
where $\bm a$ is the coefficient vector of polynomial $a$. Moreover, define the Hankel moment matrices
$$
H_k=\bigl((i+j+k)!\bigr)_{i,j=0}^m,
\qquad k=0,1,2.
$$
Then
$$
Q_k(a) = \int_0^\infty t^k a(t)^2e^{-t}dt
= \sum_{i,j=0}^m (i+j+k)!a_i a_j
=\bm a^\mathsf{T}H_k\bm a.
$$
Hence the function-space Rayleigh quotient can be written in coefficient form as
$$
\Lambda_\mu(a)
=
\frac{\bm a^\mathsf{T}
\bigl(H_2-2\mu H_1+\mu^2H_0\bigr)\bm a}
{\bm a^\mathsf{T}H_0\bm a}.
$$
Since $H_0$ is positive definite, minimizing this quotient is equivalent to a generalized eigenvalue problem. This is the matrix representation of the same first-order condition obtained above by varying the polynomial directly.
\end{remark}


\begin{proposition}\label{prop:shifted-form}
For the optimizer $a_*$ and center $\mu_*$ defined above,
\begin{equation}\label{eq:shifted-form}
  \int_0^\infty R_*(t)b(t)^2\e^{-t}\dd t\ge0,
  \qquad b\in\Pcal_m.
\end{equation}
Moreover, equality holds at $b=a_*$.
\end{proposition}
\begin{proof}
Immediate from \eqref{var:eq:eta-def}: $\eta_*\le\Lambda_{\mu_*}(b)$ for every nonzero $b\in\Pcal_m$, which rearranges to \eqref{eq:shifted-form}, with equality at $b=a_*$ since $a_*$ realizes the minimum.
\end{proof}
The conclusion of Proposition~\ref{prop:shifted-form} applies to every global optimizer of the  degree-at-most-$m$ problem, regardless of its actual degree. In particular, it remains valid if the optimizer has degree $m-1$ rather than $m$.



Set $q_* = R_* a_*$, and call $\ip{p}{q}=\int_0^\infty p(t)q(t)e^{-t}\,dt$, the weighted inner product with weight \(e^{-t}\), where $p$ and $q$ can be any polynomials.
The Euler equation gives
$$
\ip{q_*}{b} =0
\qquad
\text{for every } b\in\Pcal_m,
$$
so \(q_*\perp\Pcal_m\) in the weighted space with weight \(e^{-t}\). Since \(R_*\) is quadratic and \(a_*\in\Pcal_m\), $q_*\in\Pcal_{m+2}$.
Hence
\begin{equation}
\label{eq:q_space}
q_*\in\Pcal_{m+2}\cap\Pcal_m^\perp,
\qquad
\dim\bigl(\Pcal_{m+2}\cap\Pcal_m^\perp\bigr)=2.
\end{equation}
Hence, coefficient stationarity alone leaves a two-dimensional space of candidates for \(q_*\). Pinning down the missing direction requires further stationarity conditions.

\section{Laguerre reduction and rigidity}\label{sec:laguerre-rigidity}

The next step is to exploit the structure of the two-dimensional space $\Pcal_{m+2}\cap\Pcal_m^\perp$. For the weight $\e^{-t}$ on $(0,\infty)$, the Laguerre basis describes this space naturally, since the coefficient first-order condition forces all lower-degree Laguerre modes to vanish, leaving only two possible modes. The remaining freedom is then resolved by the damping-rate stationarity condition developed in the next step.


The ordinary Laguerre polynomials $L_n(t)=\sum_{j=0}^n(-1)^j\binom{n}{j}t^j/j!$ have exact degree $n$, with leading coefficient $(-1)^n/n!$, and satisfy the orthogonality relation
\begin{equation*}
    \ip{L_j}{L_k}
    =\delta_{jk}, \qquad j,k\mathrel{{\geq}}0,
\end{equation*}
where $\delta_{jk}$ is the Kronecker delta. Since $L_0,\ldots,L_r$ are $r+1$ polynomials of distinct degrees $0,\ldots,r$, they are linearly independent, and hence an orthonormal basis of $\Pcal_r$ for every $r\ge0$.
Then, since $q_*\in\Pcal_{m+2}$, we have a unique expansion
\begin{equation}\label{lag:eq:q-expansion}
    q_*(t)=\sum_{k=0}^{m+2}c_kL_k(t).
\end{equation}
The Euler equation will force the first $m+1$ coefficients to vanish, as we show next.
\begin{theorem}[Laguerre reduction]\label{lag:thm:reduction}
Assume \eqref{eq:q_space}.
Then
\begin{equation}\label{lag:eq:two-modes}
    R_*(t)a_*(t)
    =c_{m+1}L_{m+1}(t)+c_{m+2}L_{m+2}(t),
\end{equation}
for uniquely determined real constants $c_{m+1}$ and $c_{m+2}$.
\end{theorem}

\begin{proof}
For $j=0,\ldots,m$, $L_j\in\Pcal_m$ is an admissible test polynomial in \eqref{var:eq:Euler-integral}, so $0=\ip{q_*}{L_j}$. Substituting \eqref{lag:eq:q-expansion} and using orthonormality gives $0=\sum_{k=0}^{m+2}c_k\ip{L_k}{L_j}=c_j$, hence $c_0=\cdots=c_m=0$. This proves \eqref{lag:eq:two-modes}, with uniqueness following from linear independence of the Laguerre basis.
\end{proof}


The Laguerre reduction still leaves a two-dimensional space of candidates, and removing one of the two modes requires a variational identity of a different kind, obtained by letting the exponential damping rate itself vary. For \(\beta>0\), \(\mu>0\), and \(0\neq a\in\Pcal_m\), define
\begin{equation*}
    \mathcal J_{\beta,\mu}(a)
    =\frac{\displaystyle
      \int_0^\infty (t-\mu)^2 a(t)^2\e^{-\beta t}\dd t}
    {\displaystyle
      \mu^2\int_0^\infty a(t)^2\e^{-\beta t}\dd t}.
\end{equation*}
At \(\beta=1\), minimizing over \(\mu\) gives \(\rho(a)\), and the minimizer is \(Q_2(a)/Q_1(a)\). The parameter \(\beta\) is the exponential damping rate, and it determines the time scale of the reference exponential law. The following lemma shows that changing $\beta$ only rescales the
polynomial and the center.

\begin{lemma}\label{rig:lem:dilation}
Let \(\widetilde a(x)=a(x/\beta)\). Then \(a\mapsto\widetilde a\) is a bijection of \(\Pcal_m\), with inverse \(p(t)\mapsto p(\beta t)\), and
\begin{equation*}
    \mathcal J_{\beta,\mu}(a)
    =\mathcal J_{1,\beta\mu}(\widetilde a).
\end{equation*}
\end{lemma}
\begin{proof}
The identity follows directly from the change of variables $x=\beta t$: the
numerator and denominator acquire the same factor $\beta^{-3}$, while $a(t)$
and $\mu$ become $\widetilde a(x)$ and $\beta\mu$, respectively.
\end{proof}


In particular, varying $\beta$ changes only the time scale of the problem,
not its optimal value, which makes $\beta$ available as an additional
variational direction.

\begin{proposition}\label{rig:prop:joint-min}
The optimal value \(\rho_*\) satisfies
\begin{align*}
    \rho_*
    =\min_{\substack{\beta>0,\ \mu>0\\0\neq a\in\Pcal_m}}
      \mathcal J_{\beta,\mu}(a).
\end{align*}
Moreover, the triple \((1,\mu_*,a_*)\) is a global minimizer of the right-hand side.
\end{proposition}

\begin{proof}
By Lemma~\ref{rig:lem:dilation}, \(\mathcal J_{\beta,\mu}(a)=\mathcal J_{1,\beta\mu}(\widetilde a)\geq\rho(\widetilde a)\geq\rho_*\) for every admissible triple, while \(\mathcal J_{1,\mu_*}(a_*)=\rho(a_*)=\rho_*\). Equality at \((1,\mu_*,a_*)\) proves both assertions.
\end{proof}

The dilation identity concerns a simultaneous reparametrization of $\beta$,
the center, and the polynomial. It does not make the one-variable map $\beta{\mapsto}\mathcal J_{\beta,\mu_*}(a_*)$
constant when $a_*$ and $\mu_*$ are held fixed. Proposition~\ref{rig:prop:joint-min}
states that $(1,\mu_*,a_*)$ is a global minimizer of the enlarged
three-variable problem. Consequently, the partial derivative in the damping
direction exists and satisfies
\[
  \left.\partial_\beta\mathcal J_{\beta,\mu_*}(a_*)\right|_{\beta=1}=0,
\]
which supplies
an additional stationarity equation.

In particular, such derivative is computed by an argument identical in structure to the proof of Theorem~\ref{var:thm:first-variation}, now varying \(\beta\) in place of the polynomial direction, so we only sketch it. Differentiating the numerator and denominator of \(\mathcal J_{\beta,\mu_*}(a_*)\) under the integral sign is justified directly: for \(\beta\) in a fixed neighborhood of $1$, the differentiated integrands are bounded in absolute value by $C(1+t^{2m+3})\e^{-ct}$ for some constants $c,C>0$, which is integrable. Thus differentiation under the integral sign gives
\begin{align*}
    \partial_\beta\left[\int_0^\infty (t-\mu_*)^2a_*(t)^2\e^{-\beta t}\dd t\right]
    &=-\int_0^\infty t(t-\mu_*)^2a_*(t)^2\e^{-\beta t}\dd t,
      \\
    \partial_\beta\left[\mu_*^2\int_0^\infty a_*(t)^2\e^{-\beta t}\dd t\right]
    &=-\mu_*^2\int_0^\infty t\,a_*(t)^2\e^{-\beta t}\dd t.
\end{align*}
By Proposition~\ref{rig:prop:joint-min}, \(\beta=1\) minimizes \(\mathcal J_{\beta,\mu_*}(a_*)\), so the same quotient-rule computation as in that proof, evaluated at \(\beta=1\) with \(\mathcal J_{1,\mu_*}(a_*)=\rho_*\), gives
\begin{align*}
    0
    &=-\int_0^\infty
      t\left[(t-\mu_*)^2-\rho_*\mu_*^2\right]
      a_*(t)^2\e^{-t}\dd t.
\end{align*}
Using $\eta_*=\rho_*\mu_*^2$ and the definition of $R_*$, the vanishing derivative in $\beta$ gives, explicitly,
\begin{equation}\label{rig:eq:scale-stationarity}
    \int_0^\infty tR_*(t)a_*(t)^2\e^{-t}\dd t=0.
\end{equation}
Equivalently,
\begin{equation}\label{rig:eq:scale-inner-product}
    \ip{R_*a_*}{t a_*}=0.
\end{equation}



\begin{remark}\label{rig:rem:missing-direction}
The identity \eqref{rig:eq:scale-inner-product} has the same form as the polynomial Euler condition with $b=t a_*$. If $\deg a_*=m$, however, $t a_*\in\Pcal_{m+1}\setminus\Pcal_m$, so this direction is not covered by \eqref{var:eq:Euler-integral}. The damping-rate variation therefore supplies one additional orthogonality condition beyond those obtained from polynomial variations.
\end{remark}






In the full-degree case, this additional condition eliminates the remaining
$\Lag_{m+1}$ component. Indeed, if $\deg a_*=m$, then $t a_*$ has exact
degree $m+1$, and orthogonality of $\Lag_{m+1}$ to $\Pcal_m$ gives
$$
\ip{t a_*}{\Lag_{m+1}}
=
\frac{\lc(a_*)}{\lc(\Lag_{m+1})}
\neq 0.
$$
Here $\lc(p)$ denotes the leading coefficient of a nonzero polynomial $p$.
These observations lead to the following rigidity result, including the
possible degree-$(m-1)$ boundary case.


\begin{theorem}[One-mode Laguerre rigidity]\label{rig:thm:rigidity}
Let \(a_*\in\Pcal_m\setminus\{0\}\) be a real optimizer, and let \(R_*\) be defined by \eqref{var:eq:R-def}. Assume the coefficient stationarity condition \eqref{var:eq:Euler-integral} and the scale-stationarity condition \eqref{rig:eq:scale-stationarity}. Then $\deg a_*\in\{m-1,m\}$.
More precisely, exactly one of the following alternatives holds:
\begin{align}
    \deg a_*&=m-1,
    &R_*a_*&=c\Lag_{m+1},\label{rig:eq:lower-branch}\\
    \deg a_*&=m,
    &R_*a_*&=c\Lag_{m+2},\label{rig:eq:full-branch}
\end{align}
for some nonzero real constant \(c\).
\end{theorem}

\begin{proof}
Write \(d=\deg a_*\). Since \(R_*\) is monic of degree two and \(a_*\neq0\),
\begin{equation}\label{rig:eq:q-exact-degree}
    \deg(R_*a_*)=d+2,
    \qquad
    R_*a_*\neq0.
\end{equation}
If \(d\leq m-2\), then \(R_*a_*\in\Pcal_m\) is an admissible test polynomial in \eqref{var:eq:Euler-integral}, and taking \(b=R_*a_*\) gives \(0=\normw{R_*a_*}^2\), forcing \(R_*a_*\equiv0\) by positive definiteness, contradicting \eqref{rig:eq:q-exact-degree}. Hence \(d\geq m-1\), and since \(d\leq m\), only \(d=m-1\) and \(d=m\) remain. In the first case, \(R_*a_*\) has degree \(m+1\) and, by coefficient stationarity, is orthogonal to \(\Pcal_m\). Since the orthogonal complement of \(\Pcal_m\) inside \(\Pcal_{m+1}\) is one-dimensional, \(R_*a_*=c\Lag_{m+1}\) for some real \(c\), necessarily nonzero because \(R_*a_*\neq0\), which proves \eqref{rig:eq:lower-branch}.

In the second case \(d=m\), the polynomial \(R_*a_*\) has degree \(m+2\) and is again orthogonal to \(\Pcal_m\), so
\begin{equation}\label{rig:eq:two-mode}
    R_*a_*=c_1\Lag_{m+1}+c_2\Lag_{m+2}
\end{equation}
with \(c_2\neq0\), since \(c_2=0\) would force the right-hand side to have degree at most \(m+1\), contradicting the exact degree \(m+2\). Coefficient variation alone cannot go further, as both Laguerre components remain compatible with orthogonality to \(\Pcal_m\), so scale stationarity resolves the ambiguity. Since \(ta_*\in\Pcal_{m+1}\) and \(\Lag_{m+2}\perp\Pcal_{m+1}\), substituting \eqref{rig:eq:two-mode} into \eqref{rig:eq:scale-inner-product} gives
\[
    0=\ip{R_*a_*}{ta_*}=c_1\ip{\Lag_{m+1}}{ta_*}.
\]
The second factor is nonzero, because \(ta_*\) has exact degree \(m+1\), so \(c_1=0\) and \eqref{rig:eq:two-mode} reduces to \(R_*a_*=c_2\Lag_{m+2}\) with \(c_2\neq0\), proving \eqref{rig:eq:full-branch}.
\end{proof}

The one-mode identity has an immediate algebraic interpretation in terms of the roots of $R_*$ that are being deleted from the corresponding Laguerre polynomial. First, recall that $\rho_*\in (0,1)$. Hence, the two roots of \(R_*\) are
\begin{equation}\label{rig:eq:R-roots}
    x=\mu_*(1-\sqrt{\rho_*}),
    \qquad
    y=\mu_*(1+\sqrt{\rho_*}).
\end{equation}
They satisfy
\[
    0<x<y,
    \qquad
    R_*(t)=(t-x)(t-y).
\]

\begin{corollary}\label{rig:cor:deletion}
Under the assumptions of Theorem~\ref{rig:thm:rigidity}, let
\[
    n=
    \begin{cases}
        m+1,&\deg a_*=m-1,\\
        m+2,&\deg a_*=m.
    \end{cases}
\]
Then
\begin{equation}\label{rig:eq:Ra-Ln}
    R_*a_*=c\Lag_n
\end{equation}
for some \(c\neq0\), and the roots \(x,y\) of \(R_*\) are zeros of \(\Lag_n\). Consequently,
\begin{equation}\label{rig:eq:a-quotient}
    a_*(t)
    =c\frac{\Lag_n(t)}{(t-x)(t-y)}.
\end{equation}
\end{corollary}

\begin{proof}
Immediate from Theorem~\ref{rig:thm:rigidity}, which gives \eqref{rig:eq:Ra-Ln}. Since \((t-x)(t-y)a_*(t)=c\Lag_n(t)\) and the roots \(x,y\) are distinct by \eqref{rig:eq:R-roots}, dividing gives \eqref{rig:eq:a-quotient} and shows \(\Lag_n(x)=\Lag_n(y)=0\).
\end{proof}
First, the parameters of the quadratic have a direct interpretation. From \eqref{rig:eq:R-roots},
\begin{align*}
    \mu_* = \frac{x+y}{2},
    \qquad
    \rho_* = \left(\frac{y-x}{x+y}\right)^2,
\end{align*}
which entails that the objective value is
\begin{equation}\label{rig:eq:scv-roots}
    \Vcal(a_*)
    =\frac{\rho_*}{1-\rho_*}
    =\frac{(y-x)^2}{4xy},
\end{equation}
from target function \eqref{eq:obj_scv}.

Second, the standard Laguerre polynomial \(\Lag_n\) has \(n\) positive simple zeros. Hence, the quadratic loss polynomial removes two of those zeros, while every remaining zero becomes a zero of the half-polynomial \(a_*\). In the following section, we show that such zeros are adjacent, and we identify which pair minimizes the objective. Those conclusions require the shifted-form inequality and Gauss--Laguerre quadrature.





\section{Quadrature and the adjacent-node criterion}\label{sec:adjacency}


The preceding section reduces the optimizer to the deletion of two zeros $x<y$ from a Laguerre polynomial $\Lag_n$, with $n=m+1$ on the degree-$(m-1)$ branch and $n=m+2$ on the full-degree branch. The corresponding objective value is \eqref{rig:eq:scv-roots}.
The remaining problem is therefore to determine which pair of Laguerre zeros can occur and which such pair minimizes this quantity.
We first show, using the shifted-form inequality \eqref{eq:shifted-form} and Gauss--Laguerre quadrature, that the deleted zeros must be adjacent. This reduces the two branches to the discrete values $\Delta_{m+1}$ and $\Delta_{m+2}$, respectively. Then, we show that $\Delta_{m+2}<\Delta_{m+1}$, yielding an optimizer of exact degree $m$.




The role of quadrature is to convert the shifted-form inequality into a discrete condition at these zeros. Since the optimizer is obtained by deleting two zeros of $\Lag_n$, Gauss--Laguerre quadrature allows the relevant weighted integrals to be expressed in terms of their values at the remaining nodes. This will provide the key step in showing that the deleted zeros must be adjacent. We therefore adapt the quadrature formulas needed below; see Gautschi \cite[\S1.4.2]{gautschi2004}. In particular, the degree-$2n$ extension below is needed for the lower-degree branch, where $b\in\Pcal_{n-1}$ and therefore $R\,b^2$ may have degree $2n$.

For the next proof and the subsequent sections, write
\begin{equation}\label{adj:eq:monic-Laguerre}
 \pi_n(t)=(-1)^n n!\Lag_n(t).
\end{equation}
Then $\pi_n$ is monic, its zeros are the $n$ simple positive zeros $ 0<x_{1,n}<\cdots<x_{n,n}$,
and
\begin{equation}\label{adj:eq:monic-norm}
 \int_0^\infty \pi_n(t)^2\e^{-t}\dd t=(n!)^2.
\end{equation}

\begin{lemma}[Gauss--Laguerre quadrature]
\label{adj:lem:quadrature}
For each $n\geq1$, there are positive weights
$w_{1,n},\ldots,w_{n,n}$ such that
\begin{equation}\label{adj:eq:gauss-exact}
 \int_0^\infty p(t)\e^{-t}\dd t
 =
 \sum_{r=1}^n w_{r,n}p(x_{r,n})
\end{equation}
for every $p \in \Pcal_{2n-1}$.
Moreover, for every
$p\in\Pcal_{2n}$,
\begin{equation}\label{adj:eq:gauss-error}
 \int_0^\infty p(t)\e^{-t}\dd t
 =
 \sum_{r=1}^n w_{r,n}p(x_{r,n})
 +(n!)^2[t^{2n}]p.
\end{equation}
Here $[t^{2n}]p$ denotes the coefficient of $t^{2n}$ in $p$, which is zero whenever $\deg p<2n$.
Finally, the weights are
\begin{equation}\label{adj:eq:weight-derivative}
 w_{r,n}
 =
 \frac{1}{x_{r,n}\Lag_n'(x_{r,n})^2},
 \qquad r=1,\ldots,n.
\end{equation}
\end{lemma}

\begin{proof}
The degree-$(2n-1)$ exactness and positivity of the weights follow from
Gautschi \cite[Theorems~1.45--1.46]{gautschi2004}, while the Gaussian
remainder formula is given in
\cite[Corollary to Theorem~1.48, Eq.~(1.4.14)]{gautschi2004}.
For the monic Laguerre polynomial
$\pi_n$ in \eqref{adj:eq:monic-Laguerre}, the remainder constant is $(n!)^2$ in \eqref{adj:eq:monic-norm}.
Since
\[
\frac{p^{(2n)}}{(2n)!}=[t^{2n}]p,
\qquad p\in\Pcal_{2n},
\]
the standard remainder formula gives \eqref{adj:eq:gauss-error}.
Finally, the weight identity \eqref{adj:eq:weight-derivative} is the
Gauss--Laguerre specialization of the general Christoffel-number formula
\cite[Theorem~1.47, Eq.~(1.4.9)]{gautschi2004}.
\end{proof}








We now use quadrature to turn the shifted-form inequality into a condition on the deleted Laguerre zeros. Fix $n\ge2$ and two zeros $x_{j,n}<x_{k,n}$ of $\Lag_n$, and set
$$
R_{j,k}^{(n)}(t)=
(t-x_{j,n})(t-x_{k,n}),
$$
together with
$$
\mathfrak q_{j,k}^{(n)}(b)
=
\int_0^\infty
R_{j,k}^{(n)}(t)b(t)^2\e^{-t}\dd t.
$$
For each pair $j<k$, $R_{j,k}^{(n)}$ is the monic quadratic that would coincide with $R_*$ if $x_{j,n}$ and $x_{k,n}$ were the two deleted zeros.
Applying Gauss--Laguerre quadrature to $\mathfrak q_{j,k}^{(n)}(b)$ expresses its sign through the values of $R_{j,k}^{(n)}$ at the remaining Laguerre nodes, thereby reducing admissibility of the pair $(j,k)$ to a discrete sign condition.


\begin{theorem}[Adjacent-node criterion]\label{adj:thm:adjacent}
Let $1\leq j<k\leq n$.
\begin{enumerate}[label=\textup{(\roman*)}]
\item The quadratic form $\mathfrak q_{j,k}^{(n)}$ is nonnegative on $\Pcal_{n-2}$ if and only if $k=j+1$.
\item The same equivalence holds on $\Pcal_{n-1}$.
\item In the adjacent case, the nullspace on either space is one-dimensional and is spanned by
\begin{equation}\label{adj:eq:null-vector}
 a_{j,n}(t)
 =\frac{\Lag_n(t)}
 {(t-x_{j,n})(t-x_{j+1,n})}.
\end{equation}
The quotient in \eqref{adj:eq:null-vector} is understood as the
polynomial obtained after removing the two displayed linear factors from
$\Lag_n$.
\end{enumerate}
\end{theorem}

\begin{proof}
Assume first that $k=j+1$. Let $b\in\Pcal_{n-2}$. The polynomial $R_{j,j+1}^{(n)}b^2$ has degree at most $2+2(n-2)=2n-2$, so the exact quadrature formula \eqref{adj:eq:gauss-exact} applies:
\begin{equation}\label{adj:eq:adjacent-sum}
 \mathfrak q_{j,j+1}^{(n)}(b)
 =\sum_{r=1}^n w_{r,n}
 R_{j,j+1}^{(n)}(x_{r,n})b(x_{r,n})^2.
\end{equation}
The quadratic $R_{j,j+1}^{(n)}$ is negative only on the open interval $(x_{j,n},x_{j+1,n})$, and because the two zeros are adjacent, no quadrature node lies in that interval. Hence $R_{j,j+1}^{(n)}(x_{r,n})\geq0$ for $1\leq r\leq n$, with equality precisely for $r=j$ and $r=j+1$. Since every weight and every square in \eqref{adj:eq:adjacent-sum} is nonnegative, $\mathfrak q_{j,j+1}^{(n)}(b)\geq0$.

Suppose equality holds. Then $b$ must vanish at every node except possibly $x_{j,n}$ and $x_{j+1,n}$, and there are $n-2$ such nodes. Since $\deg b\leq n-2$, either $b=0$ or
\[
 b(t)=C\prod_{\substack{1\leq r\leq n\\r\neq j,j+1}}
 (t-x_{r,n})
\]
for some nonzero constant $C$. This product is proportional to the quotient in \eqref{adj:eq:null-vector}. Thus the nullspace on $\Pcal_{n-2}$ is exactly the stated one-dimensional span.

Now let $b\in\Pcal_{n-1}$ and write $b(t)=b_{n-1}t^{n-1}+\cdots+b_0$, where $b_{n-1}$ denotes the fixed top coordinate and may be zero. For $p(t)=R_{j,j+1}^{(n)}(t)b(t)^2$, one has $[t^{2n}]p=b_{n-1}^2$, so Equation \eqref{adj:eq:gauss-error} gives
\[
 \mathfrak q_{j,j+1}^{(n)}(b)
 =\sum_{r=1}^n w_{r,n}
 R_{j,j+1}^{(n)}(x_{r,n})b(x_{r,n})^2
 +(n!)^2b_{n-1}^2.
\]
Both terms are nonnegative. If equality holds, then $b_{n-1}=0$, hence $b\in\Pcal_{n-2}$, and the nullspace description already proved applies. This proves all positive statements.

Conversely, suppose that $k\geq j+2$. Choose an index $r$ with $j<r<k$. Then $x_{r,n}$ lies strictly between the roots of $R_{j,k}^{(n)}$, so $R_{j,k}^{(n)}(x_{r,n})<0$. Define the isolating polynomial
\begin{equation*}
 b_r(t)=
 \prod_{\substack{1\leq s\leq n\\s\neq j,r}}
 (t-x_{s,n}),
\end{equation*}
which has degree $n-2$, hence belongs to both $\Pcal_{n-2}$ and $\Pcal_{n-1}$. At the quadrature nodes, $b_r(x_{s,n})=0$ unless $s=j$ or $s=r$, and the contribution at $s=j$ is killed by $R_{j,k}^{(n)}(x_{j,n})=0$. Thus exact quadrature yields $\mathfrak q_{j,k}^{(n)}(b_r)=w_{r,n}R_{j,k}^{(n)}(x_{r,n})b_r(x_{r,n})^2<0$. The form is therefore not nonnegative even on $\Pcal_{n-2}$. This proves the converse and completes the proof.
\end{proof}




The one-mode Laguerre rigidity theorem, Theorem~\ref{rig:thm:rigidity}, identifies the roots of $R_*$
 as two Laguerre zeros, while Theorem~\ref{adj:thm:adjacent} shows that global nonnegativity of the shifted form forces those zeros to be adjacent. Applying this criterion to the two possible optimizer degrees in Theorem~\ref{rig:thm:rigidity} gives the following branchwise characterization.

\begin{proposition}\label{adj:prop:adjacency-branches}
Let $a_*$ be an optimizer of the closed problem.
\begin{enumerate}[label=\textup{(\roman*)}]
\item If $\deg a_*=m-1$, then there exists $j\in\{1,\ldots,m\}$ such that
\begin{align*}
 R_*(t)=(t-x_{j,m+1})(t-x_{j+1,m+1})
\end{align*}
and
\begin{align*}
 a_*(t)=C\frac{\Lag_{m+1}(t)}
 {(t-x_{j,m+1})(t-x_{j+1,m+1})}.
\end{align*}
\item If $\deg a_*=m$, then there exists $j\in\{1,\ldots,m+1\}$ such that
\begin{align*}
 R_*(t)=(t-x_{j,m+2})(t-x_{j+1,m+2})
\end{align*}
and
\begin{align*}
 a_*(t)=C\frac{\Lag_{m+2}(t)}
 {(t-x_{j,m+2})(t-x_{j+1,m+2})}.
\end{align*}
\end{enumerate}
\end{proposition}

\begin{proof}

In the lower branch, put $n=m+1$. Then $\Pcal_m=\Pcal_{n-1}$. The
nonnegative shifted Rayleigh form from Proposition~\ref{prop:shifted-form} is
available for this optimizer even though $\deg a_*=m-1$. It says that the
quadratic form determined by the two roots of $R_*$ is nonnegative on all of
$\Pcal_{n-1}$. Theorem~\ref{rig:thm:rigidity} places those roots among the
zeros of $\Lag_n$, and part (ii) of Theorem~\ref{adj:thm:adjacent} forces them
to be adjacent and gives the quotient form.


In the full branch, put $n=m+2$. Then $\Pcal_m=\Pcal_{n-2}$. Apply part (i) of Theorem~\ref{adj:thm:adjacent} in exactly the same way.
\end{proof}

Proposition~\ref{adj:prop:adjacency-branches} shows that, in either branch, the optimizer is obtained from the appropriate Laguerre polynomial by deleting an adjacent pair of zeros: $\Lag_{m+1}$ in the degree-$(m-1)$ branch, $\Lag_{m+2}$ in the degree-$m$ branch. Adjacency is now proved, but both degrees remain possible, since nothing so far compares the best achievable value on one branch against the best achievable value on the other. In the next section we compare them. 


\section{Exact degree and the optimal kernel}\label{sec:degree-selection}

We begin by evaluating the objective associated with deleting two zeros of a Laguerre polynomial. The resulting expressions depend only on the deleted pair and will then be specialized to the adjacent pairs identified in the previous section.

Let $n\geq2$, let $x<y$ be two distinct zeros of $\Lag_n$, and set
\begin{equation}\label{adj:eq:quotient-general}
 {a_n}(t)=\frac{\Lag_n(t)}{R(t)},
 \qquad
 R(t)=(t-x)(t-y).
\end{equation}

\begin{lemma}\label{adj:lem:moments}

For the quotient \eqref{adj:eq:quotient-general},
\begin{equation}\label{adj:eq:Q0-pair}
 Q_0({a_n})=\frac{x+y}{xy(y-x)^2},
\end{equation}
and
\begin{equation}\label{adj:eq:moment-ratios}
 \frac{Q_1({a_n})}{Q_0({a_n})}=\frac{2xy}{x+y},
 \qquad
 \frac{Q_2({a_n})}{Q_1({a_n})}=\frac{x+y}{2}.
\end{equation}
Consequently,
\begin{align*}
 \rho({a_n})=\frac{(y-x)^2}{(x+y)^2},
\end{align*}
and
\begin{equation}\label{adj:eq:V-pair}
 \Vcal({a_n})=\frac{(y-x)^2}{4xy}.
\end{equation}

\end{lemma}

\begin{proof}
Because $a_n^2$ has degree $2n-4$, the polynomial $t^ka_n(t)^2$ has degree $2n-4+k$ for each $k=0,1,2$, so Gauss--Laguerre quadrature at the $n$ zeros of $\Lag_n$ (\Cref{adj:lem:quadrature}) is exact for all three simultaneously, since $2n-4+k\le 2n-1$. After the two linear factors corresponding to $x$ and $y$ are removed, the quotient $a_n$ still contains every other linear factor of $\Lag_n$, so it vanishes at every Laguerre quadrature node except the two deleted nodes $x$ and $y$; only these two summands survive in each of the three quadrature sums
\[
 Q_k(a_n)=\int_0^\infty t^ka_n(t)^2\e^{-t}\dd t
 =x^kw_xa_n(x)^2+y^kw_ya_n(y)^2,
 \qquad k=0,1,2,
\]
where $w_x,w_y$ denote the quadrature weights at $x,y$.

At the two removable singularities,
\[
 a_n(x)=\frac{\Lag_n'(x)}{x-y},
 \qquad
 a_n(y)=\frac{\Lag_n'(y)}{y-x}.
\]
Together with \eqref{adj:eq:weight-derivative}, these identities give
\[
 w_xa_n(x)^2=\frac{1}{x(y-x)^2},
 \qquad
 w_ya_n(y)^2=\frac{1}{y(y-x)^2}.
\]
Substitution in the three quadrature sums yields

$$
Q_0(a_n) = \frac{x+y}{xy(y-x)^2}, \qquad Q_1(a_n) = \frac{2}{(y-x)^2}, \qquad Q_2(a_n) = \frac{x+y}{(y-x)^2}.
$$

Equations \eqref{adj:eq:Q0-pair} and \eqref{adj:eq:moment-ratios} follow. Finally,
\[
 \rho(a_n)=1-\frac{Q_1(a_n)^2}{Q_0(a_n)Q_2(a_n)}
 =\frac{(y-x)^2}{(x+y)^2},
\]
and \eqref{adj:eq:V-pair} follows from the relation $\Vcal(a_n)=\rho(a_n)/(1-\rho(a_n))$.
\end{proof}




For adjacent zeros, define
\begin{equation}\label{adj:eq:delta}
 {\gamma_{j,n}}
 =\frac{(x_{j+1,n}-x_{j,n})^2}
 {4x_{j,n}x_{j+1,n}},
 \qquad 1\leq j<n,
\end{equation}
and
$$
\Delta_n=\min_{1\le j<n}{\gamma_{j,n}}.
$$
By Proposition~\ref{adj:prop:adjacency-branches} and Lemma~\ref{adj:lem:moments}, every optimizer on the degree-$(m-1)$ branch has value ${\gamma_{j,m+1}}$ for some $j$, while every optimizer on the degree-$m$ branch has value ${\gamma_{j,m+2}}$ for some $j$. Consequently, the two branch optima are bounded below by $\Delta_{m+1}$ and $\Delta_{m+2}$, respectively. These bounds are attained: for every adjacent pair of zeros of $\Lag_n$, Theorem~\ref{adj:thm:adjacent} gives a nonnegative shifted form whose nullspace contains the corresponding quotient, while Lemma~\ref{adj:lem:moments} gives its objective value ${\gamma_{j,n}}$. Hence, the optimal values on the two branches are exactly $\Delta_{m+1}$ and $\Delta_{m+2}$, and therefore
\begin{equation}\label{adj:eq:closed-value}
 \mathfrak v_m^{\le}=\min\{\Delta_{m+1},\Delta_{m+2}\}{,}
\end{equation}
Here the superscript $\le$ denotes optimization over the degree-at-most-$m$ problem, allowing $\deg a\le m$ rather than requiring exact degree $m$.

To conclude $R_*a_*=c\Lag_{m+2}$, it remains to prove that $\Delta_{m+2}$ is always strictly smaller than $\Delta_{m+1}$, which requires comparing the smallest zero-gaps of consecutive Laguerre degrees. Define
\begin{align*}
 z_{j,n}=\log x_{j,n},
 \qquad
 \ell_n=\min_{1\leq j<n}(z_{j+1,n}-z_{j,n}).
\end{align*}
For $0<x<y$,
\begin{align*}
 \frac{(y-x)^2}{4xy}
 =\frac14\left(\sqrt{\frac yx}-\sqrt{\frac xy}\right)^2\notag
 =\sinh^2\left(\frac12\log\frac yx\right).
\end{align*}
Consequently,
\begin{equation}\label{adj:eq:Delta-ell}
 \Delta_n=\sinh^2(\ell_n/2).
\end{equation}
Since $u\mapsto\sinh^2(u/2)$ is strictly increasing on $(0,\infty)$, we have
$\Delta_{n+1}<\Delta_n$ if and only if $\ell_{n+1}<\ell_n$.
We therefore compare the smallest spacings of the Laguerre zeros on the logarithmic scale, using the Laguerre differential equation transformed to logarithmic coordinates. The standard Laguerre equation is
\begin{equation}\label{adj:eq:Lag-ODE}
 t\Lag_n''(t)+(1-t)\Lag_n'(t)+n\Lag_n(t)=0.
\end{equation}

\begin{lemma}\label{adj:lem:log-normal}
Define
\begin{equation}\label{adj:eq:w-def}
 w_n(s)=\exp\left(-\frac{\e^s}{2}\right)\Lag_n(\e^s),
 \qquad s\in\R.
\end{equation}
Then, the zeros of $w_n$ are $z_{1,n},\ldots,z_{n,n}$ and
\begin{equation}\label{adj:eq:w-ODE}
 w_n''(s)+W_n(s)w_n(s)=0,
 \qquad
 W_n(s)=\left(n+\frac12\right)\e^s-\frac{\e^{2s}}4.
\end{equation}
Moreover,
\begin{equation}\label{adj:eq:W-max}
 W_n(s)\leq\left(n+\frac12\right)^2,
\end{equation}
with equality only when $\e^s=2n+1$.
\end{lemma}
\begin{proof}
Set $t=\e^s$ and $u_n(s)=\Lag_n(\e^s)$, so $u_n'=t\Lag_n'(t)$ and $u_n''=t\Lag_n'(t)+t^2\Lag_n''(t)$. Multiplying \eqref{adj:eq:Lag-ODE} by $t$ and substituting gives
\begin{equation}\label{adj:eq:u-ODE}
 u_n''(s)-t u_n'(s)+nt u_n(s)=0.
\end{equation}
Writing $u_n=\e^{t/2}w_n$ and using $t'=t$,
\[
 u_n'=\e^{t/2}\left(w_n'+\frac t2w_n\right),
 \qquad
 u_n''=\e^{t/2}\left(w_n''+t w_n'+\left(\frac{t^2}{4}+\frac t2\right)w_n\right),
\]
and substituting into \eqref{adj:eq:u-ODE} cancels the first-derivative terms, leaving $w_n''+\bigl((n+\tfrac12)t-t^2/4\bigr)w_n=0$, which is \eqref{adj:eq:w-ODE} since $t=\e^s$. The exponential factor in \eqref{adj:eq:w-def} is strictly positive, so it does not change the zeros. Finally, with $k=n+1/2$,
\[
 W_n(s)=k\e^s-\frac{\e^{2s}}4=k^2-\frac14(\e^s-2k)^2,
\]
which proves \eqref{adj:eq:W-max}.
\end{proof}

The normal form \eqref{adj:eq:w-ODE} together with the bound
\eqref{adj:eq:W-max} gives a universal lower bound on the logarithmic zero
spacings. The coefficient $W_n$ controls the local oscillation rate of
solutions to
$$
w_n''+W_n(s)w_n=0.
$$
Indeed, for the constant-coefficient equation $y''+k^2y=0$, consecutive zeros
are spaced by exactly $\pi/k$. Thus, an upper bound on $W_n$ prevents the zeros
of $w_n$ from becoming arbitrarily close, and the following lemma makes this precise

\begin{lemma}\label{adj:lem:gap-lower}
For every $n\geq2$,
\begin{equation}\label{adj:eq:ell-lower}
 \ell_n\geq\frac{\pi}{n+1/2}.
\end{equation}
Consequently,
\begin{equation}\label{adj:eq:Delta-lower}
 \Delta_n\geq
 \sinh^2\left(\frac{\pi}{2n+1}\right).
\end{equation}
\end{lemma}

\begin{proof}
Let $a<b$ be consecutive zeros of $w_n$. Multiplying \eqref{adj:eq:w-ODE} by $w_n$, integrating over $[a,b]$, and using $w_n(a)=w_n(b)=0$ gives $\int_a^b w_n'(s)^2\dd s=\int_a^b W_n(s)w_n(s)^2\dd s$. The sharp Wirtinger inequality for a function vanishing at both endpoints (see, e.g., Hardy--Littlewood--P\'olya \cite[Ch.~VII, \S7.7, Eq.~(7.7.1)]{hardylittlewoodpolya1952}) gives
\[
 \int_a^b w_n'(s)^2\dd s
 \geq\frac{\pi^2}{(b-a)^2}
 \int_a^b w_n(s)^2\dd s,
\]
while the same identity also bounds the left side above by $\bigl(\max_{[a,b]}W_n\bigr)\int_a^b w_n^2\dd s$. Since $w_n$ is nonzero between consecutive zeros, its squared integral is positive, so comparing the two estimates gives
\begin{equation}\label{adj:eq:intervalwise-gap}
 \frac{\pi^2}{(b-a)^2}
 \leq \max_{s\in[a,b]}W_n(s)
 \leq\left(n+\frac12\right)^2,
\end{equation}
hence $b-a\geq\pi/(n+1/2)$. Taking the minimum over all consecutive zero intervals proves \eqref{adj:eq:ell-lower}, and \eqref{adj:eq:Delta-lower} follows from \eqref{adj:eq:Delta-ell}.
\end{proof}

The preceding lemma gives the lower bound
in \eqref{adj:eq:ell-lower}. To prove $\ell_{n+1}<\ell_n$, it therefore suffices to exhibit a consecutive degree-$(n+1)$ logarithmic gap smaller than $\pi/(n+1/2)$. The next lemma does so for $n\geq5$.

\begin{lemma}\label{adj:lem:large-comparison}
For every $n\geq5$,
\begin{equation}\label{adj:eq:strict-large}
 \ell_{n+1}<\frac{\pi}{n+1/2}\leq\ell_n.
\end{equation}
\end{lemma}

\begin{proof}
Put $k=n+\frac12$. For degree $n+1$, we have $W_{n+1}(s)=(k+1)\e^s-\e^{2s}/4$. Setting $t=\e^s$, the inequality $W_{n+1}(s)>k^2$ is equivalent to $t^2-4(k+1)t+4k^2<0$, whose roots are $t_\pm=2(k+1\pm\sqrt{2k+1})$; thus $W_{n+1}(s)>k^2$ exactly on $I_k=(\log t_-,\log t_+)$, whose length is
\begin{equation}\label{adj:eq:I-length}
 |I_k|=\log\frac{k+1+\sqrt{2k+1}}{k+1-\sqrt{2k+1}}=2\operatorname{arctanh}\left(\frac{\sqrt{2k+1}}{k+1}\right).
\end{equation}

We claim $|I_k|>2\pi/k$ for $k\geq11/2$. Let $u(k)=\sqrt{2k+1}/(k+1)$ and $h(k)=k\operatorname{arctanh}u(k)$, so $h'(k)=\operatorname{arctanh}u(k)-1/\sqrt{2k+1}$. Since $\operatorname{arctanh}u>u$ for $0<u<1$ and $u(k)>1/\sqrt{2k+1}$, we get $h'(k)>0$. At $k=11/2$,
\[
 h(11/2)>\frac{11}{2}\left[\frac{4\sqrt3}{13}+\frac13\left(\frac{4\sqrt3}{13}\right)^3\right]=\frac{4070\sqrt3}{2197}>\frac{22}{7}>\pi,
\]
the middle inequality being exact after clearing denominators and squaring: $3\cdot28490^2-48334^2=98864744>0$. Since $h$ is increasing, $h(k)>\pi$ for every $k\geq11/2$, which by \eqref{adj:eq:I-length} is the claim \eqref{adj:eq:I-long}:
\begin{equation}\label{adj:eq:I-long}
 |I_k|>\frac{2\pi}{k},
 \qquad k\geq\frac{11}{2}.
\end{equation}

Choose $[A,B]\subset I_k$ with $B-A>2\pi/k$. Let $w=w_{n+1}$ and, since a nontrivial solution of a second-order linear equation cannot have $w,w'$ vanish simultaneously, define a continuous Pr\"ufer angle by $\tan\theta(s)=kw(s)/w'(s)$, equivalently $w'=r\cos\theta$, $kw=r\sin\theta$ for $r=(w'^2+k^2w^2)^{1/2}>0$. Then $w(s)=0$ if and only if $\theta(s)\in\pi\mathbb Z$, so consecutive
zeros of $w$ are exactly the points at which $\theta$ attains consecutive
multiples of $\pi$, separated by the increase in $s$ needed for $\theta$ to
grow by $\pi$. This is the standard Pr\"ufer-angle device (see, e.g.,
\cite{teschl2012}). Differentiating and using $w''=-W_{n+1}w$,
\begin{equation*}
 \theta'(s)=k\frac{w'(s)^2+W_{n+1}(s)w(s)^2}{w'(s)^2+k^2w(s)^2},
\end{equation*}
so $\theta'(s)\geq k$ on $I_k$, with equality only where $w(s)=0$, which is a measure-zero set, since $w$ is a nontrivial solution of a second-order linear equation. For comparison, the constant-coefficient equation $y''+k^2y=0$ has $\theta(s)=ks$ for $y=\sin(ks)$, i.e.\ zeros spaced exactly $\pi/k$ apart; here $W_{n+1}>k^2$ on $I_k$ makes $\theta$ advance strictly faster. Integrating over $[A,B]$ therefore gives $\theta(B)-\theta(A)>k(B-A)>2\pi$.

Since $\theta$ is continuous and strictly increasing with
$\theta(B)-\theta(A)>2\pi$, the interval $(\theta(A),\theta(B))$ contains two
consecutive multiples of $\pi$, and the points $c<d$ at which $\theta$
attains these two values are consecutive zeros of $w$ with
$\theta(d)-\theta(c)=\pi$ exactly. Integrating $\theta'\geq k$ over $[c,d]$ gives $\pi>k(d-c)$, i.e.\ $d-c<\pi/k$. Since $w_{n+1}(s)=\e^{-\e^s/2}L_{n+1}(\e^s)$ vanishes exactly where $s=\log x_{r,n+1}$ for a zero $x_{r,n+1}$ of $L_{n+1}$, we have $c=\log x_{r,n+1}$, $d=\log x_{r+1,n+1}$ for some $r$, so $d-c$ is one of the logarithmic gaps minimized by $\ell_{n+1}$; hence
\[
 \ell_{n+1}\leq d-c<\frac\pi k.
\]
Combined with $\pi/k\leq\ell_n$ from Lemma~\ref{adj:lem:gap-lower}, this proves \eqref{adj:eq:strict-large}.
\end{proof}


Lemma~\ref{adj:lem:large-comparison} begins at $n=5$. The remaining comparisons $n=2,3,4$ are obtained by direct numerical evaluation in the next lemma.

\begin{lemma}\label{adj:lem:small-comparison}
One has
\begin{equation}\label{adj:eq:small-comparisons}
 \Delta_3<\Delta_2,
 \qquad
 \Delta_4<\Delta_3,
 \qquad
 \Delta_5<\Delta_4.
\end{equation}
\end{lemma}
\begin{proof}
Direct numerical evaluation of the Laguerre zeros and the relative-gap
score in \eqref{adj:eq:delta} gives
\[
 \Delta_2=1,\qquad
 \Delta_3\approx0.2765825308,\qquad
 \Delta_4\approx0.1384532326,\qquad
 \Delta_5\approx0.0861276525.
\]
These values immediately give
$\Delta_3<\Delta_2$, $\Delta_4<\Delta_3$, and
$\Delta_5<\Delta_4$, proving \eqref{adj:eq:small-comparisons}.
\end{proof}

Combining the finite-degree comparisons in Lemma~\ref{adj:lem:small-comparison} ($n=2,3,4$) with the large-degree comparison in Lemma~\ref{adj:lem:large-comparison} ($n\geq5$), together with \eqref{adj:eq:Delta-ell}, gives strict monotonicity of $\Delta_n$ at every degree:
\begin{equation}\label{adj:eq:Delta-monotone}
 \Delta_{n+1}<\Delta_n,
 \qquad n\geq2.
\end{equation}
This is also an independent spacing result for Laguerre zeros in its own right: the minimum logarithmic spacing decreases strictly with the degree, a comparison that ordinary interlacing does not by itself imply (compare the standard oscillation framework in \cite{szego}). This is exactly what is needed to rule out the lower-degree branch and select the full-degree optimizer. We can now exclude the boundary branch and state the optimal kernel explicitly.

\begin{theorem}[Optimal repeated-pole rational kernel]\label{adj:thm:final}
Let $a_*$ be a global optimizer of the closed degree-at-most-$m$ polynomial-square problem, and let $R_*$ be defined by \eqref{var:eq:R-def}. Then
\begin{align*}
 \deg a_*=m,
 \qquad
 R_*(t)a_*(t)=c\Lag_{m+2}(t),
\end{align*}
for some nonzero real constant $c$. Moreover, there is an index $j\in\{1,\ldots,m+1\}$ such that
\begin{align*}
 R_*(t)
 =(t-x_{j,m+2})(t-x_{j+1,m+2}),
\end{align*}
and
\begin{equation}\label{adj:eq:final-a}
 a_*(t)
 =C\frac{\Lag_{m+2}(t)}
 {(t-x_{j,m+2})(t-x_{j+1,m+2})}.
\end{equation}
The deleted pair minimizes the relative adjacent-gap score, and this common value is the optimum $\mathfrak v_m$ itself:
\begin{equation}\label{adj:eq:final-value}
\mathfrak v_m
 =\Vcal(a_*)
 =\Delta_{m+2}
 =\min_{1\leq r\leq m+1}
 \frac{(x_{r+1,m+2}-x_{r,m+2})^2}
 {4x_{r,m+2}x_{r+1,m+2}}.
\end{equation}
At the endpoint $m=0$, this reads $\mathfrak v_0=\Delta_2=1$. The unique half-polynomial is constant and the mean-one kernel is the unit exponential density.

Writing $n=m+2$, $x=x_{j,n}$, and $y=x_{j+1,n}$ for the minimizing pair, so that $a_*(t)=C\,\Lag_n(t)/\bigl((t-x)(t-y)\bigr)$, the normalized unit-damping density is
\begin{equation}\label{main:eq:unit-damping-opt}
 \kappa_{1,a_*}(z)
 =\frac{xy(y-x)^2}{x+y}\e^{-z}
 \left[\frac{\Lag_n(z)}{(z-x)(z-y)}\right]^2,
\end{equation}
and
\begin{equation}\label{main:eq:unit-moments}
 \E[T_{a_*}]=\frac{2xy}{x+y},
 \qquad
 \E[T_{a_*}^2]=xy.
\end{equation}
Choosing $\lambda_{m,j}=2xy/(x+y)$ produces the optimal mean-one kernel
\begin{align*}
 \kappa_{m,j}^*(u)
 =\lambda_{m,j}\frac{xy(y-x)^2}{x+y}
 \e^{-\lambda_{m,j}u}
 \left[
 \frac{\Lag_n(\lambda_{m,j}u)}
 {(\lambda_{m,j}u-x)(\lambda_{m,j}u-y)}
 \right]^2,
\end{align*}
with removable singularities, for which
\begin{align*}
 \int_0^\infty\kappa_{m,j}^*(u)\dd u=1,
 \qquad
 \E[U_{m,j}^*]=1,
 \qquad
 \Var(U_{m,j}^*)={\gamma_{j,n}}=\Delta_n.
\end{align*}
No other kernel in the exact-degree repeated-pole square-polynomial class has smaller variance, and every optimal kernel arises this way, up to the irrelevant sign and magnitude of the half-polynomial, from one minimizing adjacent pair.
\end{theorem}

\begin{proof}
By \eqref{adj:eq:closed-value}, the best value in the closed problem is
\[
 \mathfrak v_m^{\leq}=\min\{\Delta_{m+1},\Delta_{m+2}\}.
\]
Equation \eqref{adj:eq:Delta-monotone}, applied with $n=m+1$, gives
\[
 \Delta_{m+2}<\Delta_{m+1}.
\]
Therefore
\begin{equation}\label{adj:eq:closed-selected}
 \mathfrak v_m^{\leq}=\Delta_{m+2}.
\end{equation}
Every degree-$(m-1)$ optimizer belongs to the $\Lag_{m+1}$ branch and, by Proposition~\ref{adj:prop:adjacency-branches} and Lemma~\ref{adj:lem:moments}, has value at least $\Delta_{m+1}$. This is strictly larger than \eqref{adj:eq:closed-selected}. Hence no global optimizer can have degree $m-1$.

Theorem~\ref{rig:thm:rigidity} leaves only the full-degree alternative, proving
\[
 \deg a_*=m,
 \qquad
 R_*a_*=c\Lag_{m+2}.
\]
Proposition~\ref{adj:prop:adjacency-branches} then forces the two roots of $R_*$ to be an adjacent pair of zeros of $\Lag_{m+2}$ and gives the quotient representation \eqref{adj:eq:final-a}. Finally, because the global value is $\Delta_{m+2}$, the selected adjacent pair must attain the minimum in \eqref{adj:eq:final-value}, which proves $\mathfrak v_m=\Delta_{m+2}$.

For $m\geq1$, this adjacent-deletion form and the value $\Delta_{m+2}$ give, via Equation \eqref{adj:eq:Q0-pair}, the absolute
normalization in \eqref{main:eq:unit-damping-opt}, while
\eqref{adj:eq:moment-ratios} gives \eqref{main:eq:unit-moments}. Rescaling by
$\lambda_{m,j}=\E[T_{a_*}]$ produces a mean-one density without changing the
squared coefficient of variation, and \eqref{adj:eq:V-pair} gives its
variance. For $m=0$, $n=2$, the unique quotient is constant; using
$x+y=4$, $xy=2$, and $(y-x)^2=8$ reduces the formulas to
$\kappa_{1,a_*}(z)=\e^{-z}$, $\lambda=1$, and $\mathfrak v_0=1=\Delta_2$.
\end{proof}

Expanding the squared optimal half-polynomial as in
Remark~\ref{aw:rem:confluent} writes every optimal kernel as a signed
combination of common-rate Erlang densities. Any calculation linear in the
randomizing law can therefore be carried out for the required Erlang orders
and recombined with the optimal signed weights. This is the
same linearity exploited by the concentrated matrix-exponential inversion
method, which evaluates numerical Laplace-inversion formulas through an
analogous signed combination
\cite{horvathtalyigastelek2018,horvathetal2020nilt,almousaetal2022cme}.

The theorem reduces the original optimization over a continuous family of polynomial-square kernels to a finite problem on the zeros of a single Laguerre polynomial. In particular, the degree relaxation is exact, the optimizer has full degree $m$, and every optimal kernel is obtained by deleting an adjacent pair of zeros of $\Lag_{m+2}$ that minimizes the relative-gap score. Hence, both the optimal value and the corresponding mean-one kernel are determined explicitly by the zero geometry of $\Lag_{m+2}$.


\section{Constructive computation}\label{sec:computation}
We turn briefly to computational aspects of the construction. Most of what follows is standard, but it is worth stating explicitly in the present setting. 
In particular, take $n=m+2$. By \cref{adj:thm:final}, the nonlinear projective
optimization is exactly the finite comparison
\begin{equation}\label{num:eq:exact-input}
 \mathfrak v_m
 =\min_{1\le j<n}
 \frac{(x_{j+1,n}-x_{j,n})^2}{4x_{j,n}x_{j+1,n}}.
\end{equation}
This section gives a direct construction of the zeros and the minimizing
adjacent pair.

Recall the monic Laguerre polynomial $\pi_n(t)=(-1)^n n!\Lag_n(t)$ from
\eqref{adj:eq:monic-Laguerre}. This satisfies
\begin{equation}\label{num:eq:monic-recurrence}
 \pi_n(t)
 =
 \bigl(t-(2n-1)\bigr)\pi_{n-1}(t)
 -(n-1)^2\pi_{n-2}(t),
\end{equation}
with $\pi_0=1$ and $\pi_1=t-1$. Moreover, let
\begin{align*}
 \mathsf J_n=
 \begin{pmatrix}
  1 & 1 \\
  1 & 3 & 2 \\
    & 2 & 5 & \ddots \\
    &   & \ddots & \ddots & n-1\\
    &   &        & n-1 & 2n-1
 \end{pmatrix},
\end{align*}
be the Jacobi matrix associated with the monic Laguerre recurrence.
\begin{proposition}\label{num:prop:Jacobi}
For every $n\ge1$,
\[
 \det(tI_n-\mathsf J_n)=\pi_n(t).
\]
The eigenvalues of $\mathsf J_n$ are therefore
$x_{1,n},\ldots,x_{n,n}$. They are positive and simple.
\end{proposition}

\begin{proof}
The leading principal determinants of $tI_n-\mathsf J_n$ satisfy the
continuant recurrence \eqref{num:eq:monic-recurrence} with the same initial
values as $\pi_n$. This proves the characteristic-polynomial identity.
Moreover,
\[
 \mathsf J_n=B_nB_n^{\mathsf T},
\]
where $B_n$ is lower bidiagonal with
\[
 (B_n)_{kk}=\sqrt{k},
 \qquad
 (B_n)_{k,k-1}=\sqrt{k-1}.
\]
Thus $\mathsf J_n$ is positive definite. Its off-diagonal entries are
nonzero, so it is an irreducible real symmetric tridiagonal matrix and has
simple spectrum.
\end{proof}

Since the zeros are exactly the eigenvalues of the symmetric tridiagonal
Jacobi matrix $\mathsf J_n$, they can be computed with
specialized tridiagonal eigensolvers, avoiding direct root finding for the
Laguerre polynomial. The optimizer can therefore be constructed from this spectrum by the procedure in Algorithm~\ref{alg:optimizer}.

\begin{algorithm}[htbp]
\caption{Optimizer via the Jacobi spectrum}
\label{alg:optimizer}
\begin{algorithmic}[1]
\State $n \gets m+2$
\State Form the diagonal $d_k=2k-1$ and off-diagonal $e_k=k$ of $\mathsf J_n$
\State Compute the ordered eigenvalues $x_1<\cdots<x_n$ with a symmetric tridiagonal eigensolver
\For{$j = 1, \ldots, n-1$}
    \State ${\gamma_{j}} \gets \sinh^2\!\left[\dfrac12\log\!\left(\dfrac{x_{j+1}}{x_j}\right)\right]$
\EndFor
\State \Return any $j_* \in \argmin_{1\leq j<n}{\gamma_{j}}$
\end{algorithmic}
\end{algorithm}
Standard algorithms for real symmetric tridiagonal matrices
compute this spectrum; see \cite{lapack1999}.
The same tridiagonal eigenvalue problem underlies the Golub--Welsch algorithm
for Gauss-Laguerre quadrature nodes and weights \cite{golub1969}, treated in detail,
together with the numerical construction of orthogonal polynomials more
generally, in \cite{gautschi2004}.

Table~\ref{tab:finite-values} lists the resulting values of
$\mathfrak v_m$ for $m=0,\ldots,10$ against the Erlang benchmark $1/N$. Each row requires only forming $\mathsf J_{m+2}$, computing its
ordered eigenvalues, and scanning the adjacent relative gaps. No nonlinear
optimization is needed, and the results already show the strict improvement
over the Erlang benchmark from $m=1$ onward. The ten-digit display is retained
to facilitate comparison with the repeated-pole values reported by
M\'esz\'aros and Telek \cite{meszarostelek2022}. As noted in the
introduction, the admissible class here is larger than their real-rooted
parametrization, yet the optimizer is itself real-rooted, so the two minima
coincide and the values above supply the global certificate their numerical
study did not provide.

\begin{table}[h!]
\centering
\caption{Finite-order values obtained directly from the Jacobi
spectrum. The pair index is one-based, $N=2m+1$ is the rational order, and
the Erlang benchmark has the same order. Values are rounded to ten decimal
places.}
\label{tab:finite-values}
\small

\begin{tabular}{@{}rrrrrr@{}}
\toprule
$m$ & $N$ & $j_*$ & $\mathfrak v_m$ & $1/N$ & $N^2\mathfrak v_m$\\
\midrule
$0$  & $1$  & $1$  & $1.0000000000$ & $1.0000000000$ & $1.0000000000$\\
$1$  & $3$  & $2$  & $0.2765825308$ & $0.3333333333$ & $2.4892427770$\\
$2$  & $5$  & $3$  & $0.1384532326$ & $0.2000000000$ & $3.4613308143$\\
$3$  & $7$  & $4$  & $0.0861276525$ & $0.1428571429$ & $4.2202549738$\\
$4$  & $9$  & $5$  & $0.0600486248$ & $0.1111111111$ & $4.8639386117$\\
$5$  & $11$ & $6$  & $0.0449172542$ & $0.0909090909$ & $5.4349877526$\\
$6$  & $13$ & $7$  & $0.0352403354$ & $0.0769230769$ & $5.9556166832$\\
$7$  & $15$ & $7$  & $0.0283976276$ & $0.0666666667$ & $6.3894662039$\\
$8$  & $17$ & $8$  & $0.0228210879$ & $0.0588235294$ & $6.5952943947$\\
$9$  & $19$ & $9$  & $0.0188503929$ & $0.0526315789$ & $6.8049918389$\\
$10$ & $21$ & $10$ & $0.0159086737$ & $0.0476190476$ & $7.0157251231$\\
\bottomrule
\end{tabular}

\end{table}

Beyond the finite-order computation of the optimizer, the same Jacobi spectra
also reveal a clear pattern in the location of the minimizing deleted pair.
As $m\to\infty$, its midpoint divided by $m$ appears to approach $2$, while
the gap between the two deleted zeros approaches $2\pi$.
Figure~\ref{fig:finite-localization} illustrates this behavior at finite $m$. The next section proves these limits and uses them to determine
the large-order behavior of the optimal value.

\begin{figure}[h]
\centering
\begin{minipage}[t]{0.485\textwidth}
\centering
\includegraphics[width=\linewidth]{figures/localization_midpoint_m1000.pdf}
\end{minipage}\hfill
\begin{minipage}[t]{0.485\textwidth}
\centering
\includegraphics[width=\linewidth]{figures/localization_gap_m1000.pdf}
\end{minipage}
\caption{Finite-order localization diagnostics obtained from the same Jacobi
spectra as \cref{tab:finite-values}, computed for $1\leq m\leq1000$. The left
panel shows the midpoint of a minimizing deleted pair divided by $m$; the
dashed line is its limit $2$. The right panel shows the corresponding
absolute gap; the dashed line is its limit $2\pi$.
}
\label{fig:finite-localization}
\end{figure}



\section{Large-order asymptotics and localization}\label{sec:asymptotics}

Let $j_n$ be any index attaining the minimum in the definition of $\Delta_n$.
By \eqref{adj:eq:Delta-ell}, this pair also realizes the minimum logarithmic
spacing $\ell_n$, while Lemma~\ref{adj:lem:gap-lower} gives
$\ell_n\geq\pi/(n+1/2)$. The sharp asymptotic constant is obtained from the
limiting distribution of the normalized Laguerre zeros.

Define the empirical measure $\nu_n=n^{-1}\sum_{j=1}^n\delta_{x_{j,n}/n}$. The relevant result is the following specialization of Kornyik and Michaletzky \cite[Theorem~1]{kornyikmichaletzky2017}: for the zeros of $L_p^{(\alpha_p)}$ with $\alpha_p/p\to c>-1$, the normalized empirical measure has limiting support $x_\pm=(\sqrt{c+1}\pm1)^2$ and, when $c\geq0$, limiting density $\sqrt{(x_+-x)(x-x_-)}/(2\pi x)$ on $[x_-,x_+]$. Here, $p=n$, $\alpha_p=0$, $c=0$, so $x_-=0$, $x_+=4$, and
\begin{align*}
 \rho_{\MP}(x)
 =
 \frac1{2\pi}\sqrt{\frac{4-x}{x}},
 \qquad 0<x<4,
\end{align*}
giving $\nu_n\Rightarrow\nu_{\MP}$
and, since $\nu_{\MP}(\partial I)=0$ for every closed interval $I\subset(0,4)$,
\begin{equation}\label{asy:eq:continuity}
 \nu_n(I)\longrightarrow\nu_{\MP}(I).
\end{equation}

\begin{theorem}[Sharp asymptotics and localization]\label{thm:large-order}
For every choice of minimizing indices $j_n$,
\begin{align}
 n^2\Delta_n&\longrightarrow\frac{\pi^2}{4},
 \label{asy:eq:value}\\
 n\ell_n&\longrightarrow\pi,
 \label{asy:eq:log}\\
 \frac{x_{j_n,n}}n,\ 
 \frac{x_{j_n+1,n}}n&\longrightarrow2,
 \label{asy:eq:location}\\
 x_{j_n+1,n}-x_{j_n,n}&\longrightarrow2\pi,
 \label{asy:eq:absolute}\\
 \frac{j_n}{n}&\longrightarrow\frac12+\frac1\pi.
 \label{asy:eq:index}
\end{align}
\end{theorem}

\begin{proof}
From \eqref{adj:eq:ell-lower} and \eqref{adj:eq:Delta-ell}, $\liminf_n n^2\Delta_n\ge\pi^2/4$. For the upper bound, fix $\varepsilon\in(0,1)$, put $I_\varepsilon=[2-\varepsilon,2+\varepsilon]$, $\mu_\varepsilon=\nu_{\MP}(I_\varepsilon)$, and let $N_n(\varepsilon)=\#\{j:x_{j,n}/n\in I_\varepsilon\}$, so that $N_n(\varepsilon)/n\to\mu_\varepsilon$ by \eqref{asy:eq:continuity}. Since $\mu_\varepsilon>0$, eventually $N_n(\varepsilon)\ge2$; writing $r_n,s_n$ for the first and last indices in $I_\varepsilon$, $\sum_{j=r_n}^{s_n-1}(x_{j+1,n}-x_{j,n})=x_{s_n,n}-x_{r_n,n}\le2\varepsilon n$, so some adjacent pair $k_n$ satisfies $x_{k_n+1,n}-x_{k_n,n}\le2\varepsilon n/(N_n(\varepsilon)-1)$. Both zeros are at least $n(2-\varepsilon)$, so
\[
 \Delta_n
 \le
 \frac{\varepsilon^2}
 {(2-\varepsilon)^2\{N_n(\varepsilon)-1\}^2}
 \quad\Longrightarrow\quad
 \limsup_n n^2\Delta_n
 \le
 \frac{\varepsilon^2}{(2-\varepsilon)^2\mu_\varepsilon^2}.
\]
Since $\rho_{\MP}$ is continuous at $2$ with $\rho_{\MP}(2)=1/(2\pi)$, $\mu_\varepsilon=\varepsilon/\pi+o(\varepsilon)$, and letting $\varepsilon\downarrow0$ gives $\limsup_n n^2\Delta_n\le\pi^2/4$, proving \eqref{asy:eq:value}. Since $\ell_n=2\operatorname{arsinh}\sqrt{\Delta_n}$ and $\operatorname{arsinh}u/u\to1$, this also gives \eqref{asy:eq:log}.

For the localization statements, set $x_n=x_{j_n,n}$, $y_n=x_{j_n+1,n}$, and recall the logarithmic normal form $w_n''+W_nw_n=0$ with $W_n(s)=(n+\tfrac12)\e^s-\e^{2s}/4$. Since larger $W_n$ permits shorter zero spacing 
(\Cref{adj:eq:intervalwise-gap})
near equality in the Poincar\'e--Wirtinger bound can occur only where $W_n$ is close to its maximum; applying the intervalwise estimate \eqref{adj:eq:intervalwise-gap} on $[\log x_n,\log y_n]$, where $W_n$ is continuous and hence attains its maximum, gives a point $s_n$ in this interval with $W_n(s_n)\ge\pi^2/\ell_n^2$. Setting $t_n=\e^{s_n}\in[x_n,y_n]$, \eqref{adj:eq:w-ODE} divided by $n^2$ gives, using $n\ell_n\to\pi$,
\[
 \left(1+\frac1{2n}\right)\frac{t_n}{n}
 -
 \frac14\left(\frac{t_n}{n}\right)^2
 \ge1+o(1),
\]
while completing the square on the same expression shows it equals $(1+1/2n)^2-\tfrac14(t_n/n-2-1/n)^2\le(1+1/2n)^2=1+o(1)$. Squeezing these two bounds forces the subtracted square to vanish asymptotically, so $t_n/n\to2$; since $x_n\le t_n\le y_n$ and $y_n/x_n=\e^{\ell_n}\to1$, also $x_n/t_n\to1$ and $y_n/t_n\to1$, proving \eqref{asy:eq:location}. Then $y_n-x_n=(x_n/n)(n\ell_n)(\e^{\ell_n}-1)/\ell_n\to2\pi$, proving \eqref{asy:eq:absolute}.

Finally, for $a\in(0,4)$ let $M_n(a)=\#\{j:x_{j,n}/n\le a\}$. Localization gives, for every fixed $\varepsilon\in(0,2)$ and large $n$, $M_n(2-\varepsilon)<j_n\le M_n(2+\varepsilon)$; dividing by $n$, by the weak convergence $\nu_n\Rightarrow\nu_{\MP}$, and letting $\varepsilon\downarrow0$ gives $j_n/n\to F_{\MP}(2)=\int_0^2\rho_{\MP}$. Substituting $x=4\sin^2\theta$, $F_{\MP}(2)=\tfrac4\pi\int_0^{\pi/4}\cos^2\theta\,\dd\theta=\tfrac12+\tfrac1\pi$, proving \eqref{asy:eq:index}.
\end{proof}

The theorem is stated in terms of the Laguerre degree $n$. We first translate these limits to the original polynomial degree $m$ and rational order $N=2m+1$.

\begin{corollary}\label{asy:cor:m}
Let $\mathfrak v_m$ be the optimum in
\eqref{num:eq:exact-input}. Then
\begin{align}
 \mathfrak v_m
 &=
 \frac{\pi^2}{4m^2}+o(m^{-2})\notag\\
 \label{eq:conv_spacing}
 (2m+1)^2\mathfrak v_m
 &\longrightarrow\pi^2.
\end{align}
For every minimizing pair among the zeros of $\Lag_{m+2}$,
\begin{align*}
 \frac{x_{j_m,m+2}}{m+2},\
 \frac{x_{j_m+1,m+2}}{m+2}
 &\longrightarrow2,\\
 x_{j_m+1,m+2}-x_{j_m,m+2}
 &\longrightarrow2\pi,\\
 \frac{j_m}{m+2}
 &\longrightarrow\frac12+\frac1\pi.
\end{align*}
\end{corollary}

\begin{proof}
Set $n=m+2$ in \cref{thm:large-order} and use
\eqref{num:eq:exact-input}. Since $m/(m+2)\to1$, the first limit follows.
The exact rational order is $2m+1$, whose ratio to $m$ tends to $2$.
The remaining statements are direct translations of
\eqref{asy:eq:location}--\eqref{asy:eq:index} (in finite-$n$ terms the
leading location is $2n=2m+4$, asymptotically equivalent to $2m$).
\end{proof}

Consequently, the optimal value decays quadratically in the rational order, with sharp constant $\pi^2$, while the minimizing deleted pair localizes near $2(m+2)$ with asymptotic separation $2\pi$. We next translate these geometric asymptotics into the behavior of the corresponding optimal randomizer.

\begin{corollary}\label{asy:cor:moments}
Let $T_m^*$ have the optimal unit-damping density

\[
\kappa_{1,a_*}(u)
=
\frac{\e^{-u}a_*(u)^2}{Q_0(a_*)},
\qquad u\ge0.
\]
Then
\[
 \frac{\E[T_m^*]}{m+2}\longrightarrow2,
 \qquad
 \Var(T_m^*)\longrightarrow\pi^2.
\]
After rescaling to mean one, the optimal randomizers satisfy $\E[(U_m^*-1)^2]\longrightarrow0$, and hence concentrate at $1$ in mean square and in probability.
\end{corollary}

\begin{proof}
Let $x_m<y_m$ be the deleted pair. By \eqref{num:eq:mean}--\eqref{num:eq:second} with $\lambda=1$, and Lemma~\ref{adj:lem:moments},
\[
 \E[T_m^*]=\frac{2x_my_m}{x_m+y_m},
 \qquad
 \Var(T_m^*)
 =
 \frac{x_my_m(y_m-x_m)^2}{(x_m+y_m)^2}.
\]
Applying \cref{asy:cor:m}, the first ratio tends to $2(m+2)$, the dimensionless factor in the variance tends to $1/4$, and $(y_m-x_m)^2\to4\pi^2$. Finally, $\E[(U_m^*-1)^2]=\Var(U_m^*)=\mathfrak v_m\to0$, and Chebyshev's inequality gives $\mathbb P\{|U_m^*-1|>\varepsilon\}\le\mathfrak v_m/\varepsilon^2\to0$ for every $\varepsilon>0$, the asserted convergence in probability.
\end{proof}

Hence, the optimal unit-damping kernels retain an asymptotically finite absolute spread while their mean grows linearly with the order. After normalization to mean one, this spread vanishes, and the optimal randomizers converge to the deterministic target. Together with \eqref{eq:conv_spacing},
this gives both the sharp concentration rate and its probabilistic interpretation.


\section{Discussion}\label{sec:discussion}
Theorem~\ref{adj:thm:final} solves a precise approximation problem. For the positive
operator $\Aop_m^*h(t_0)=\E[h(t_0U_m^*)]$, \cref{rand:prop:error-bound} gives
\[
 \sup_{\lVert h''\rVert_\infty\le1}
 |\Aop_m^*h(t_0)-h(t_0)|
 =\frac{t_0^2}{2}\mathfrak v_m
 =\left(\frac{\pi^2}{8m^2}+o(m^{-2})\right)t_0^2.
\]
Since $\mathfrak v_m$ is by construction the minimal variance in
the class, this is also the minimal worst-case error attainable by any
admissible kernel. However, the construction does more than bound this error uniformly. By \cref{asy:cor:moments}, the mean-one randomizer $U_m^*$ itself concentrates at $1$ in mean square and in probability as $m\to\infty$. Thus, the randomized evaluation $t_0U_m^*$ is not merely close to the deterministic point $t_0$ in a worst-case sense at each finite order. It is a genuinely concentrating family of random approximations to that point, sharpening toward it as $m$ grows. 

At rational order $N=2m+1$, the Erlang member of the class has variance
$1/N$, whereas
\[
 N^2\mathfrak v_m\longrightarrow\pi^2.
\]
The square-polynomial matrix-exponential class therefore
concentrates at rate $N^{-2}$, a full order in $N$ better than Erlang's
$N^{-1}$. The finite values in
\cref{tab:finite-values} show the improvement already at small orders. This
comparison is a concentration statement, not a universal replacement theorem
for Erlangization: phase-type recursions may not extend unchanged to the
optimal matrix-exponential kernels.

The fixed-pole approximation literature optimizes different quantities.
Andersson's concentrated-pole problem and the related fixed-real-pole results
\cite{andersson1981,anderssonganelius1977,borwein1983} concern global
approximation errors for rational functions. The present admissible set is constrained by time-domain
nonnegativity and a square-polynomial factor, while its criterion is the worst-case error above. Consequently, the optimal damping and asymptotic
constants answer a different question, even though the repeated-real-pole
architecture is closely related.

Within the broader CME-R programme, two comparisons with M\'esz\'aros and
Telek \cite{meszarostelek2022} are worth keeping separate. The present
class is a proper subclass of their full CME-R problem, since all
transform poles are forced to coalesce at one negative real location, but
an enlargement of the real-rooted repeated-pole family used in their
numerical study, since $a$ need not be real-rooted. As noted in the
introduction, the resulting optimizer is nonetheless real-rooted, so the
two minima coincide and the present result supplies a global analytic
certificate for the optimum they observed numerically.

The broader CME and CME-R constructions
\cite{horvathsafartelekzambo2016,horvathhorvathtelek2020,almousatelek2021,meszarostelek2022}
allow larger parameterizations and are generally obtained numerically. At
some finite orders these broader parameterizations report smaller SCVs, but
the reported values come from numerical optimization rather than global
certificates for their full admissible classes. The optimization-free
complex-pole family of Battagliola and Peralta
\cite{battagliolaperalta2026} addresses another part of the design landscape.
Its squared coefficient of variation has order $\log^3(N)/N^2$, which decays more slowly than the $N^{-2}$ rate proved here, although its construction is more
explicit and avoids the Laguerre geometry and quadrature machinery used in this work.

The constructive mechanism proceeds in four stages: projective minimization over the closed polynomial class, one-mode Laguerre rigidity, adjacent-zero deletion, and the Jacobi spectral construction. Natural extensions include the full nonnegative-polynomial cone, several distinct real damping rates, and conjugate pairs of complex damping rates. The one-mode rigidity and adjacency arguments are specific to the present subclass and
should not be assumed to survive unchanged.

Taken together, these results give a complete solution of the repeated-real-pole
square-polynomial problem, both at finite order and in the large-order limit.
The optimizer is determined by the zero geometry of a single Laguerre
polynomial, can be recovered from a symmetric tridiagonal eigenvalue problem,
and yields the sharp concentration law $N^2\mathfrak v_m\longrightarrow\pi^2$.

\begin{thebibliography}{99}

\bibitem{abatewhitt2006}
J.~Abate and W.~Whitt,
A unified framework for numerically inverting Laplace transforms,
\emph{INFORMS J. Comput.} \textbf{18} (2006), 408--421.
\url{https://doi.org/10.1287/ijoc.1050.0137}

\bibitem{akargursoyhorvathtelek2021}
N.~Akar, \"O.~G\"ursoy, G.~Horv\'ath, and M.~Telek,
Transient and first passage time distributions of first- and second-order
multi-regime Markov fluid queues via ME-fication,
\emph{Methodol. Comput. Appl. Probab.} \textbf{23} (2021), 1257--1283.
\url{https://doi.org/10.1007/s11009-020-09812-y}

\bibitem{aldousshepp1987}
D.~Aldous and L.~Shepp,
The least variable phase type distribution is Erlang,
\emph{Comm. Statist. Stochastic Models} \textbf{3} (1987), 467--473.
\url{https://doi.org/10.1080/15326348708807067}

\bibitem{almousaetal2022cme}
S.~Al-Deen Almousa, G.~Horv\'ath, I.~Horv\'ath, A.~M\'esz\'aros, and M.~Telek,
The CME method: Efficient numerical inverse Laplace transformation with
concentrated matrix exponential distribution,
\emph{SIGMETRICS Perform. Eval. Rev.} \textbf{49} (2022), 29--34.
\url{https://doi.org/10.1145/3543146.3543155}

\bibitem{almousatelek2021}
S.~Al-Deen Almousa and M.~Telek,
Enhanced optimization of high order concentrated matrix-exponential
distributions,
\emph{Ann. Math. Inform.} \textbf{53} (2021), 5--19.
\url{https://doi.org/10.33039/ami.2021.02.001}

\bibitem{lapack1999}

E.~Anderson et al.,
\emph{LAPACK Users' Guide}, 3rd ed.,
Society for Industrial and Applied Mathematics, Philadelphia, 1999.
\url{https://doi.org/10.1137/1.9780898719604}


\bibitem{andersson1981}

J.-E.~Andersson,
Approximation of $\e^{-x}$ by rational functions with concentrated negative
poles,
\emph{J. Approx. Theory} \textbf{32} (1981), 85--95.
\url{https://doi.org/10.1016/0021-9045(81)90106-4}


\bibitem{anderssonganelius1977}

J.-E.~Andersson and T.~Ganelius,
The degree of approximation by rational functions with fixed poles,
\emph{Math. Z.} \textbf{153} (1977), 161--166.


\bibitem{asmussenavramusabel2002}
S.~Asmussen, F.~Avram, and M.~Usabel,
Erlangian approximations for finite-horizon ruin probabilities,
\emph{ASTIN Bull.} \textbf{32} (2002), 267--281.
\url{https://doi.org/10.2143/AST.32.2.1029}

\bibitem{bakergravesmorris}
G.~A.~Baker, Jr. and P.~Graves-Morris,
\emph{Pad\'e Approximants}, 2nd ed.,
Cambridge University Press, Cambridge, 1996.

\bibitem{battagliolaperalta2026}
M.~L.~Battagliola and O.~Peralta,
Optimization-free concentrated matrix-exponentials,
arXiv:2604.26304 (preprint).


\bibitem{bladtnielsen2017}
M.~Bladt and B.~F.~Nielsen,
\emph{Matrix-Exponential Distributions in Applied Probability},
Probability Theory and Stochastic Modelling, vol.~81,
Springer, New York, 2017.
\url{https://doi.org/10.1007/978-1-4939-7049-0}


\bibitem{bladtnielsenperalta2019}
M.~Bladt, B.~F.~Nielsen, and O.~Peralta,
Parisian types of ruin probabilities for a class of dependent risk-reserve
processes,
\emph{Scand. Actuar. J.} \textbf{2019} (2019), 32--61.
\url{https://doi.org/10.1080/03461238.2018.1483420}

\bibitem{borwein1983}

P.~B.~Borwein,
Rational approximations with real poles to $\e^{-x}$ and $x^n$,
\emph{J. Approx. Theory} \textbf{38} (1983), 279--283.
\url{https://doi.org/10.1016/0021-9045(83)90134-X}


\bibitem{boslevenbergortega2021}
L.~Bos, N.~Levenberg, and J.~Ortega-Cerd\`a,
Optimal polynomial prediction measures and extremal polynomial growth,
\emph{Constr. Approx.} \textbf{54} (2021), 431--453.
\url{https://doi.org/10.1007/s00365-020-09522-1}

\bibitem{bouchardelkarouitouzi2005}
B.~Bouchard, N.~El~Karoui, and N.~Touzi,
Maturity randomization for stochastic control problems,
\emph{Ann. Appl. Probab.} \textbf{15} (2005), 2575--2605.
\url{https://doi.org/10.1214/105051605000000593}

\bibitem{bultheeletal1999}
A.~Bultheel, P.~Gonz\'alez-Vera, E.~Hendriksen, and O.~Nj{\aa}stad,
\emph{Orthogonal Rational Functions},
Cambridge University Press, Cambridge, 1999.

\bibitem{carr1998}
P.~Carr,
Randomization and the American put,
\emph{Rev. Financ. Stud.} \textbf{11} (1998), 597--626.
\url{https://doi.org/10.1093/rfs/11.3.597}

\bibitem{gautschi2004}
W.~Gautschi,
\emph{Orthogonal Polynomials: Computation and Approximation},
Oxford University Press, Oxford, 2004.

\bibitem{golub1969}
G.~H.~Golub and J.~H.~Welsch,
Calculation of Gauss quadrature rules,
\emph{Math. Comp.} \textbf{23} (1969), 221--230.

\bibitem{hardylittlewoodpolya1952}
G.~H.~Hardy, J.~E.~Littlewood, and G.~P\'olya,
\emph{Inequalities}, 2nd ed.,
Cambridge University Press, Cambridge, 1952.

\bibitem{horvathhorvathtelek2020}
G.~Horv\'ath, I.~Horv\'ath, and M.~Telek,
High order concentrated matrix-exponential distributions,
\emph{Stochastic Models} \textbf{36} (2020), 176--192.
\url{https://doi.org/10.1080/15326349.2019.1702058}

\bibitem{horvathetal2020nilt}
I.~Horv\'ath, G.~Horv\'ath, S.~Al-Deen Almousa, and M.~Telek,
Numerical inverse Laplace transformation using concentrated matrix exponential
distributions,
\emph{Performance Evaluation} \textbf{137} (2020), 102067.
\url{https://doi.org/10.1016/j.peva.2019.102067}

\bibitem{horvathsafartelekzambo2016}
I.~Horv\'ath, O.~S\'af\'ar, M.~Telek, and B.~Z\'amb\'o,
Concentrated matrix exponential distributions,
in \emph{Computer Performance Engineering},
Lecture Notes in Computer Science, vol.~9951,
Springer, Cham, 2016, pp.~18--31.
\url{https://doi.org/10.1007/978-3-319-46433-6_2}

\bibitem{horvathtalyigastelek2018}
I.~Horv\'ath, Z.~Talyig\'as, and M.~Telek,
An optimal inverse Laplace transform method without positive and negative
overshoot---an integral based interpretation,
\emph{Electron. Notes Theor. Comput. Sci.} \textbf{337} (2018), 87--104.
\url{https://doi.org/10.1016/j.entcs.2018.03.035}

\bibitem{kornyikmichaletzky2017}
M.~Kornyik and G.~Michaletzky,
On the moments of roots of Laguerre-polynomials and the
Marchenko--Pastur law,
\emph{Ann. Univ. Sci. Budapest. Sect. Comput.} \textbf{46} (2017),
137--151.
\url{https://doi.org/10.71352/ac.46.137}

\bibitem{meszarostelek2022}
A.~M\'esz\'aros and M.~Telek,
Concentrated matrix exponential distributions with real eigenvalues,
\emph{Probab. Engrg. Inform. Sci.} \textbf{36} (2022), 1171--1187.
\url{https://doi.org/10.1017/S0269964821000309}


\bibitem{peralta2023}
O.~Peralta,
A Markov jump process associated with the matrix-exponential distribution,
\emph{J. Appl. Probab.} \textbf{60} (2023), 1--13.
\url{https://doi.org/10.1017/jpr.2022.25}


\bibitem{prevostrivoal2007}
M.~Pr\'evost and T.~Rivoal,
Remainder Pad\'e approximants for the exponential function,
\emph{Constr. Approx.} \textbf{25} (2007), 109--123.
\url{https://doi.org/10.1007/s00365-006-0635-6}

\bibitem{szego}
G.~Szeg\H{o},
\emph{Orthogonal Polynomials}, 4th ed.,
American Mathematical Society Colloquium Publications, vol.~23,
American Mathematical Society, Providence, RI, 1975.

\bibitem{teschl2012}
G.~Teschl,
\emph{Ordinary Differential Equations and Dynamical Systems},
Graduate Studies in Mathematics, vol.~140,
American Mathematical Society, Providence, RI, 2012.

\end{thebibliography}

\end{document}
